% ApuntesRedes_Parte2.tex — Parte 2: Nivel Senior
\documentclass[12pt, a4paper]{book}

% Encoding & Language
\usepackage[utf8]{inputenc}
\usepackage[T1]{fontenc}
\usepackage[spanish, es-tabla]{babel}

% Emojis (requiere LuaLaTeX)
\usepackage{emoji}
\setemojifont{Segoe UI Emoji}

% Iconos extra
\usepackage{fontawesome5}

% Layout
\usepackage[top=2.5cm, bottom=2.5cm, left=3cm, right=3cm]{geometry}
\usepackage{hyperref}
\hypersetup{
    colorlinks=true,
    linkcolor=blue!70!black,
    urlcolor=blue!70!black,
    pdftitle={Apuntes de Redes de Datos --- Parte 2},
    pdfauthor={Aldo Zepeda}
}

% Cajas visuales
\usepackage[most]{tcolorbox}
\tcbuselibrary{skins, breakable}

% Cajas personalizadas
\newtcolorbox{concepto}[1][]{
    title={\emoji{brain} Concepto},
    colback=blue!8!white, colframe=blue!50!black,
    fonttitle=\bfseries, breakable, #1
}
\newtcolorbox{ejemplo}[1][]{
    title={\emoji{package} Ejemplo real},
    colback=green!8!white, colframe=green!50!black,
    fonttitle=\bfseries, breakable, #1
}
\newtcolorbox{cuidado}[1][]{
    title={\emoji{warning} Cuidado},
    colback=orange!15!white, colframe=orange!70!black,
    fonttitle=\bfseries, breakable, #1
}
\newtcolorbox{moderno}[1][]{
    title={\emoji{rocket} Tecnología Moderna},
    colback=violet!8!white, colframe=violet!60!black,
    fonttitle=\bfseries, breakable, #1
}
\newtcolorbox{codigo}[1][]{
    title={\emoji{laptop} Configuración / Código},
    colback=gray!10!white, colframe=gray!60!black,
    fonttitle=\bfseries, breakable, #1
}

% Código fuente
\usepackage{listings}
\lstset{
    basicstyle=\ttfamily\small,
    backgroundcolor=\color{gray!10},
    frame=single,
    breaklines=true
}

% Diagramas
\usepackage{tikz}
\usetikzlibrary{shapes, arrows.meta, positioning, calc, decorations.pathreplacing}

% Matemáticas
\usepackage{amsmath}

% Colores
\usepackage{xcolor}

% Tablas
\usepackage{booktabs}
\usepackage{multirow}
\usepackage{array}

% Gráficos
\usepackage{graphicx}
\graphicspath{{assets/}}

% Macro para emojis en headings (evita MakeUppercase bug)
\newcommand{\eh}[1]{\texorpdfstring{\emoji{#1}}{}}

%% ============================================================
\begin{document}

% Evitar que \MakeUppercase del book class rompa emojis en running headers
\renewcommand{\chaptermark}[1]{\markboth{\chaptername\ \thechapter.\ #1}{}}
\renewcommand{\sectionmark}[1]{\markright{\thesection.\ #1}}

% Portada
\begin{titlepage}
    \centering
    \vspace*{2cm}
    {\Huge\bfseries \emoji{rocket} Apuntes de Redes de Datos\par}
    \vspace{0.5cm}
    {\Large\bfseries Parte 2 --- Nivel Senior\par}
    \vspace{1cm}
    {\large SDN, Datacenter, Cloud, Kubernetes, MPLS, Multicast\par}
    \vspace{0.5cm}
    {\large Aldo Zepeda Moctezuma\par}
    \vspace{0.5cm}
    {\normalsize \today\par}
    \vspace{2cm}
    {\large\itshape ``La infraestructura que no entiendes es la infraestructura que te va a fallar a las 3am.''\\[4pt]
    --- Sabiduría de SRE\par}
\end{titlepage}

\tableofcontents
\newpage

% Capítulos
\chapter{\eh{rocket} Introducción a la Parte 2 --- Nivel Senior}

%% ── 0.1 ¿Qué cubre esta parte? ──────────────────────────────────────────────
\section{\eh{books} ¿Qué cubre esta parte?}

\begin{concepto}
La \textbf{Parte 1} de estos apuntes cubrió los \textbf{fundamentos de redes}:
el modelo OSI, TCP/IP, protocolos clave (TCP, UDP, DNS, HTTP, TLS, SSH) y
cómo fluyen los datos de una máquina a otra.

La \textbf{Parte 2} sube un nivel. Aquí vemos cómo funcionan las redes
\textbf{por dentro}, en producción real, a escala empresarial y de nube:

\begin{itemize}
    \item \textbf{Datacenters modernos:} VxLAN, EVPN, Spine-Leaf, SDN
    \item \textbf{Redes de operador (carrier):} MPLS, QoS, ingeniería de tráfico
    \item \textbf{Kubernetes:} CNI, Service Mesh, Ingress, eBPF
    \item \textbf{Cloud networking:} AWS VPC, Transit Gateway, Direct Connect
    \item \textbf{Seguridad de switch empresarial:} 802.1X, MACsec, VLAN hardening
    \item \textbf{Wi-Fi empresarial:} WPA3-Enterprise, 802.11ax (Wi-Fi 6)
    \item \textbf{Multicast:} PIM, IGMP, streaming a gran escala
\end{itemize}

\emoji{light-bulb} \textbf{Analogía:}
La Parte 1 es aprender a \textbf{manejar un auto}: sabes presionar el acelerador,
girar el volante, respetar semáforos.
La Parte 2 es aprender a \textbf{construir el motor}: sabes cómo funciona
el sistema de combustión, la transmisión, el turbocompresor --- y por qué
falla a las 3am cuando hay 50{,}000 usuarios conectados.
\end{concepto}

%% ── 0.2 El Stack del Datacenter Moderno ─────────────────────────────────────
\section{\eh{office-building} El Stack del Datacenter Moderno}

Antes de entrar a cada tecnología, necesitas ver el panorama completo.
Este es el flujo real de un paquete desde Internet hasta el contenedor
que lo procesa en un datacenter moderno (AWS, Google Cloud, un banco grande):

\begin{center}
\begin{lstlisting}
Internet
   |
   v
Edge Routers  (BGP — hablan con todo Internet)
   |
   v
Firewall / DDoS Mitigation  (filtra ataques, limita flujos)
   |
   v
Load Balancers  (L4/L7 — distribuyen tráfico entre servidores)
   |
   v
Spine Switches  (ECMP — múltiples rutas iguales, sin STP)
   |
   v
Leaf Switches  (VxLAN VTEP — virtualizan la red sobre IP)
   |
   v
Servers / Hypervisors  (Linux + Open vSwitch o eBPF)
   |
   v
Pods / Containers  (Kubernetes — la app vive aquí)
\end{lstlisting}
\end{center}

\begin{concepto}
Cada flecha del diagrama es un capítulo de esta parte:

\begin{itemize}
    \item \textbf{BGP} se estudia en Cloud Networking (cap.\ 07).
    \item \textbf{Firewall / DDoS} es parte de Seguridad empresarial (cap.\ 02).
    \item \textbf{Load Balancers + VxLAN + Spine-Leaf} son el cap.\ 04 (Datacenter).
    \item \textbf{MPLS y QoS} mueven paquetes entre datacenters (cap.\ 05).
    \item \textbf{Pods y Kubernetes} son el cap.\ 06.
\end{itemize}

\emoji{light-bulb} \textbf{Analogía:}
El datacenter es como un aeropuerto enorme. Internet es el mundo exterior.
Los Edge Routers son la \textit{aduana}. Los Spine Switches son las
\textit{pistas de aterrizaje} (no hay semáforos --- todo va en paralelo).
Los Leaf Switches son las \textit{puertas de abordaje}. Los Pods son
\textit{los asientos de cada vuelo}: el destino final del pasajero (paquete).
\end{concepto}

%% ── 0.3 Mapa de Tecnologías ─────────────────────────────────────────────────
\section{\eh{world-map} Mapa de Tecnologías}

La siguiente tabla muestra qué cubre cada capítulo, qué tecnologías
involucra y en qué entorno vive en producción real:

\begin{center}
\renewcommand{\arraystretch}{1.4}
\begin{tabular}{>{\bfseries}c p{4.5cm} p{5.5cm}}
\toprule
\textbf{Cap.} & \textbf{Tecnología} & \textbf{Dónde vive} \\
\midrule
01 & Wi-Fi Empresarial: 802.11ax, WPA3-Enterprise, RADIUS, roaming 802.11r
   & Oficinas corporativas, campus universitarios, hoteles, aeropuertos \\
02 & Switching Avanzado: VLANs, STP/RSTP, 802.1X, MACsec, port security
   & Cualquier red con switches administrables (empresa, datacenter, banco) \\
03 & SDN y OpenFlow: controladores, APIs northbound/southbound, OpenDaylight
   & Datacenters de hiperescaladores (Google, Facebook), redes universitarias \\
04 & Datacenter: VxLAN, EVPN, Spine-Leaf, ECMP, BGP-EVPN
   & AWS, Azure, GCP, Equinix, datacenters privados grandes \\
05 & MPLS y QoS: LSP, LDP, RSVP-TE, ingeniería de tráfico, DiffServ
   & Redes de operador (Telmex, AT\&T, Lumen), troncales entre datacenters \\
06 & Kubernetes Networking: CNI, Calico, Cilium/eBPF, Service Mesh, Ingress
   & Toda empresa que corre microservicios en producción \\
07 & Cloud Networking: VPC, Transit Gateway, Direct Connect, BGP con AWS
   & AWS, Azure, GCP --- el 90\% de las startups y empresas medianas \\
08 & Multicast: PIM-SM, PIM-SSM, IGMP, MDT, streaming y finanzas
   & Broadcasters (televisión IP), bolsas de valores, CDNs \\
\bottomrule
\end{tabular}
\end{center}

%% ── 0.4 Parte 1 vs Parte 2 ──────────────────────────────────────────────────
\section{\eh{balance-scale} Parte 1 vs Parte 2 --- ¿Qué cambia?}

Esta tabla compara el nivel de cada tema entre las dos partes.
No es que Parte 1 sea ``incorrecto'' --- es que Parte 2 es lo que
se usa en producción a escala:

\begin{center}
\renewcommand{\arraystretch}{1.4}
\begin{tabular}{p{3.8cm} p{4cm} p{4cm}}
\toprule
\textbf{Dimensión} & \textbf{Parte 1 (Fundamentos)} & \textbf{Parte 2 (Producción)} \\
\midrule
Protocolos de transporte
   & TCP / UDP básico, puertos, handshake
   & MPLS / VxLAN, túneles, encapsulación multi-capa \\
Seguridad
   & TLS (HTTPS), SSH, firewalls básicos
   & MACsec (capa 2), 802.1X, WPA3-Enterprise, Zero Trust \\
Wi-Fi
   & 802.11 básico, WPA2 personal, canales
   & Wi-Fi 6 (802.11ax), roaming empresarial, RADIUS/EAP \\
Enrutamiento
   & RIP, OSPF básico, tablas de rutas
   & BGP completo, MPLS-TE, ECMP en Spine-Leaf \\
Switches
   & VLANs básicas, STP para evitar loops
   & EVPN, VxLAN VTEP, MLAG, port security avanzado \\
Virtualización de red
   & No aplica (redes físicas)
   & SDN, OpenFlow, OVS, eBPF, Cilium \\
Contenedores
   & No aplica
   & CNI plugins, Kubernetes networking, Service Mesh \\
Cloud
   & No aplica
   & VPC, Transit Gateway, Direct Connect, Cloud Router \\
Escala
   & Red de laboratorio o SOHO (10--100 hosts)
   & Miles de servidores, millones de flujos simultáneos \\
\bottomrule
\end{tabular}
\end{center}

%% ── 0.5 Por qué importa ─────────────────────────────────────────────────────
\section{\eh{graduation-cap} ¿Por qué importa todo esto?}

\begin{moderno}
\textbf{El 90\% de la infraestructura empresarial moderna corre sobre:}

\begin{itemize}
    \item \textbf{Kubernetes} para orquestar contenedores
    \item \textbf{Cloud networking} (AWS VPC, Azure VNet) para conectar todo
    \item \textbf{MPLS} en las redes de operador que unen datacenters
    \item \textbf{VxLAN + EVPN} dentro del datacenter para virtualizar la red
    \item \textbf{802.1X + MACsec} para seguridad de acceso en campus empresariales
\end{itemize}

\vspace{0.5em}
\textbf{Las entrevistas de SRE / Network Engineer en empresas como Google,
Amazon, Banco de México, BBVA, Telmex y cualquier empresa con más de
500 empleados preguntan exactamente estos temas.}

\vspace{0.5em}
\emoji{brain} Si sabes explicar por qué VxLAN usa UDP port 4789,
cómo Kubernetes asigna IPs a los pods, y qué hace un LSR en MPLS ---
ya estás en el top 10\% de candidatos.

\vspace{0.5em}
\emoji{light-bulb} \textbf{Analogía final:}
Un médico general sabe que el corazón bombea sangre (Parte 1).
Un cirujano cardiovascular sabe hacer un bypass, reparar una válvula
y gestionar la UCI post-operatoria (Parte 2).
Ambos son necesarios, pero el hospital solo le confía el quirófano al segundo.
\end{moderno}

\chapter{\eh{satellite-antenna} Wireless Empresarial --- WPA3 y 802.1X}

\begin{concepto}
El Wi-Fi empresarial moderno no es el mismo router de casa con una contraseña
compartida. En entornos corporativos, aeropuertos, universidades y hospitales,
la seguridad inalámbrica se construye sobre tres pilares:
\textbf{autenticación individual}, \textbf{cifrado por sesión} y
\textbf{control de acceso por puerto}.

Este capítulo cubre \textbf{WPA3} (el protocolo de seguridad Wi-Fi más moderno)
y \textbf{802.1X} (el estándar de control de acceso a red por puerto, usado tanto
en Wi-Fi como en switches cableados).
\end{concepto}

%% ── 1.1 WPA3 ────────────────────────────────────────────────────────────────
\section{\eh{locked} WPA3 --- Wi-Fi Protected Access 3}

\begin{concepto}
\textbf{WPA2} (802.11i, 2004) usa \textbf{PSK} (Pre-Shared Key): una sola
contraseña compartida que todos los usuarios conocen. Quien sabe la contraseña
entra. Si esa contraseña se filtra, \textit{toda la red queda comprometida}.
Además, un atacante puede capturar el 4-way handshake y atacarlo offline
con diccionario (Hashcat a millones de intentos por segundo).

\textbf{WPA3} (2018) reemplaza PSK por \textbf{SAE} (Simultaneous Authentication
of Equals), también conocido como ``Dragonfly Handshake''. Cada cliente se
autentica individualmente con un protocolo de prueba de conocimiento cero.

\emoji{light-bulb} \textbf{Analogía:}
\begin{itemize}
    \item \textbf{WPA2-PSK} = llave maestra que todos los empleados tienen.
          Si un empleado la pierde (o la fotografía), cualquiera entra.
    \item \textbf{WPA3-SAE} = tarjeta inteligente individual por empleado.
          Cada tarjeta es única; perder una no compromete las demás.
          Además, la tarjeta no ``revela'' el PIN al sistema --- solo demuestra
          que lo conoce (prueba de conocimiento cero).
\end{itemize}
\end{concepto}

\begin{center}
\renewcommand{\arraystretch}{1.5}
\begin{tabular}{>{\bfseries}p{3.8cm} p{4.5cm} p{4.5cm}}
\toprule
\textbf{Característica} & \textbf{WPA2} & \textbf{WPA3} \\
\midrule
Intercambio de llaves
    & PSK (Pre-Shared Key) --- contraseña compartida
    & SAE / Dragonfly --- Diffie-Hellman con prueba de conocimiento cero \\
Forward Secrecy
    & No --- capturar el handshake permite descifrar sesiones pasadas
    & Sí --- cada sesión genera llaves efímeras únicas \\
Ataques offline
    & Vulnerable --- atacante captura el 4-way handshake y ataca con diccionario
    & Resistente --- SAE no permite ataques offline de diccionario \\
Redes abiertas
    & Sin cifrado (aeropuertos, cafeterías en texto claro)
    & OWE (Opportunistic Wireless Encryption) --- cifrado sin contraseña \\
Modo empresarial
    & WPA2-Enterprise con 802.1X/EAP
    & WPA3-Enterprise con 802.1X/EAP + Suite B (criptografía de 192 bits) \\
Estándar IEEE
    & 802.11i (2004)
    & IEEE 802.11-2020 (enmienda) \\
\bottomrule
\end{tabular}
\end{center}

%% ── 1.1.1 SAE ───────────────────────────────────────────────────────────────
\subsection{\eh{key} SAE --- Simultaneous Authentication of Equals}

\begin{concepto}
\textbf{SAE} (RFC 7664 / IEEE 802.11-2020) es un protocolo de establecimiento
de llaves basado en \textbf{Diffie-Hellman sobre curva elíptica} con una
\textbf{prueba de conocimiento cero} (\textit{zero-knowledge proof}).

\textbf{¿Qué significa prueba de conocimiento cero aquí?}
En WPA2, durante el 4-way handshake, el cliente y el AP derivan una PMK
a partir de la contraseña y la intercambian de forma que un atacante que capture
el tráfico puede verificar contraseñas offline. En SAE, el protocolo demuestra
que \textit{ambas partes conocen la misma contraseña} sin que ninguna la transmita
--- ni en texto claro ni en forma derivada atacable offline.

\vspace{0.5em}
\textbf{Forward Secrecy (Secreto hacia adelante):}
Cada sesión Wi-Fi negocia un par de llaves efímeras (de un solo uso).
Si un atacante graba todo el tráfico cifrado de hoy y mañana roba
la contraseña del AP, \textit{no puede descifrar el tráfico de hoy}
porque las llaves efímeras ya no existen. WPA2 no tiene esta propiedad.
\end{concepto}

\begin{cuidado}
\textbf{KRACK Attack (Key Reinstallation Attack, 2017) --- CVE-2017-13077}

Mathy Vanhoef descubrió que en el 4-way handshake de WPA2, un atacante
Man-in-the-Middle podía forzar la \textbf{reinstalación de llaves ya usadas}
manipulando los mensajes del handshake. Al reinstalar una llave con un
nonce repetido, el cifrado CCMP se debilita drásticamente --- en algunos
casos hasta recuperación de texto claro.

\textbf{Impacto:} afectó prácticamente todo dispositivo Wi-Fi con WPA2 del planeta.

\textbf{WPA3/SAE es inmune por diseño:} SAE no usa un 4-way handshake con
reinstalación de llaves. El protocolo Dragonfly genera llaves únicas por sesión
que no pueden ser reinstaladas, eliminando la clase de vulnerabilidad que
KRACK explotó.

\emoji{light-bulb} Parches de WPA2 existen, pero WPA3 es la solución estructural.
\end{cuidado}

%% ── 1.1.2 OWE ───────────────────────────────────────────────────────────────
\subsection{\eh{globe-with-meridians} OWE --- Opportunistic Wireless Encryption}

\begin{concepto}
\textbf{OWE} (RFC 8110) resuelve un problema clásico: las redes Wi-Fi abiertas
(aeropuertos, cafeterías, hoteles) no tienen contraseña, por lo que el tráfico
viaja en texto claro y cualquier persona con un adaptador en modo monitor
puede \textbf{sniffear pasivamente} todo lo que transmites.

\textbf{OWE} mantiene la experiencia de usuario (no pide contraseña, te conectas
directamente) pero cifra automáticamente el enlace inalámbrico mediante un
intercambio \textbf{Diffie-Hellman} entre cliente y AP:

\begin{enumerate}
    \item Cliente y AP intercambian claves públicas efímeras DH durante la asociación.
    \item Derivan una \textbf{PMK} (Pairwise Master Key) compartida.
    \item El 4-way handshake estándar procede para generar llaves de sesión.
    \item Todo el tráfico queda cifrado con AES-CCM, incluso sin contraseña.
\end{enumerate}

\emoji{warning} OWE protege contra \textbf{sniffing pasivo} (alguien escucha tu
tráfico Wi-Fi). No autentica al AP --- un Evil Twin todavía puede responder.
Para autenticación del AP, se necesita 802.1X o certificados adicionales.
\end{concepto}

%% ── 1.2 802.1X ──────────────────────────────────────────────────────────────
\section{\eh{shield} 802.1X --- Port-Based Network Access Control}

\begin{concepto}
\textbf{802.1X} (IEEE 802.1X-2020) es el estándar de \textbf{control de acceso
a red basado en puerto}. Define tres actores:

\begin{itemize}
    \item \textbf{Supplicant} (Suplicante): el cliente que quiere acceso a la red
          (laptop, teléfono, impresora). Implementado en el sistema operativo
          (wpa\_supplicant en Linux, native supplicant en Windows/macOS).
    \item \textbf{Authenticator} (Autenticador): el switch o Access Point.
          Tiene el puerto en estado \textbf{UNAUTHORIZED} (bloqueado) hasta que
          el servidor RADIUS aprueba al cliente. Solo deja pasar tráfico EAP
          (EAPoL) durante la autenticación.
    \item \textbf{Authentication Server} (Servidor de Autenticación): típicamente
          un servidor \textbf{RADIUS} (FreeRADIUS, Cisco ISE, Microsoft NPS).
          Verifica las credenciales y devuelve Access-Accept o Access-Reject.
\end{itemize}

\emoji{light-bulb} \textbf{Analogía:} el sistema de llave magnética de un hotel.
\begin{itemize}
    \item El \textbf{Supplicant} es el huésped que quiere entrar a su habitación.
    \item El \textbf{Authenticator} es la cerradura electrónica de la puerta.
          No importa que el huésped sepa el número de habitación (contraseña
          compartida): sin la tarjeta válida, la puerta no abre.
    \item El \textbf{RADIUS Server} es el sistema central del hotel que verifica
          si esa tarjeta corresponde a una reserva activa y autoriza el acceso.
\end{itemize}

Si la tarjeta no es válida (credenciales incorrectas, certificado expirado,
dispositivo no administrado), el puerto queda en UNAUTHORIZED y el cliente
solo puede acceder a una VLAN de cuarentena --- si es que eso se configura.
\end{concepto}

%% ── 1.2.1 EAP ───────────────────────────────────────────────────────────────
\subsection{\eh{puzzle-piece} EAP --- Extensible Authentication Protocol}

\textbf{EAP} (RFC 3748) es el protocolo de autenticación que viaja dentro de
802.1X. Es ``extensible'' porque define solo el \textit{marco} del intercambio
de autenticación; el mecanismo real es el \textbf{método EAP} elegido:

\begin{center}
\renewcommand{\arraystretch}{1.5}
\begin{tabular}{>{\bfseries}p{2.8cm} p{3.5cm} p{4.2cm} p{2.5cm}}
\toprule
\textbf{Método} & \textbf{Mecanismo} & \textbf{Caso de uso} & \textbf{Estado} \\
\midrule
EAP-TLS
    & Certificados X.509 mutuos (cliente y servidor)
    & Máxima seguridad; dispositivos administrados con PKI propia
    & \textcolor{green!60!black}{\textbf{Recomendado}} \\
PEAP
    & Contraseña dentro de túnel TLS (solo el servidor tiene certificado)
    & Integración con Active Directory; BYOD con credenciales de dominio
    & Aceptable \\
EAP-TTLS
    & Contraseña dentro de túnel TLS con PAP/CHAP interno
    & Linux, Android, BYOD heterogéneo; no requiere certificado en el cliente
    & Aceptable \\
EAP-MD5
    & Hash MD5 de la contraseña (sin cifrado del canal)
    & Legacy; sin protección contra ataques de diccionario offline
    & \textcolor{red}{\textbf{Obsoleto --- No usar}} \\
EAP-FAST
    & PAC (Protected Access Credential) en lugar de certificado
    & Alternativa a PEAP en entornos Cisco sin PKI completa
    & Uso limitado \\
\bottomrule
\end{tabular}
\end{center}

%% ── 1.2.2 RADIUS ────────────────────────────────────────────────────────────
\subsection{\eh{satellite} RADIUS --- Remote Authentication Dial-In User Service}

\textbf{RADIUS} (RFC 2865/2866) es el protocolo de AAA (Authentication,
Authorization, Accounting) más usado en redes empresariales. Opera sobre UDP:

\begin{lstlisting}
  Puertos RADIUS:
    UDP 1812  — Authentication (sustituye al histórico 1645)
    UDP 1813  — Accounting     (sustituye al histórico 1646)

  Mensajes principales:
    Authenticator --> RADIUS:  Access-Request
                               (contiene: User-Name, NAS-IP, EAP-Message,
                                Message-Authenticator [HMAC-MD5])

    RADIUS --> Authenticator:  Access-Accept
                               (+ atributos: Tunnel-Type=VLAN,
                                Tunnel-Medium-Type=802,
                                Tunnel-Private-Group-ID=<VLAN-ID>)

                               Access-Reject
                               (acceso denegado, puede incluir Reply-Message)

                               Access-Challenge
                               (solicita más información al cliente,
                                p.ej. segundo factor, certificado)

  Seguridad:
    - El secreto compartido (RADIUS shared secret) cifra los atributos
      sensibles (contraseñas) usando MD5.
    - Para producción: usar RADSEC (RADIUS over TLS, RFC 6614)
      en lugar de RADIUS sobre UDP plano.
\end{lstlisting}

%% ── 1.2.3 Flujo completo 802.1X ─────────────────────────────────────────────
\subsection{\eh{right-arrow-curving-up} Flujo Completo 802.1X --- Diagrama de Secuencia}

El siguiente diagrama muestra el flujo completo de autenticación 802.1X
con EAP-TLS entre los tres actores:

\begin{center}
\begin{tikzpicture}[
    msg/.style={-{Stealth}, thick},
    msgr/.style={{Stealth}-, thick},
    lbl/.style={font=\small\ttfamily, text=textlight},
    lbls/.style={font=\scriptsize\itshape, text=gray!60!white}
]

%% ── Lifelines ───────────────────────────────────────────────────────────────
\node[font=\bfseries\small, align=center, text=textlight] (sup) at (0, 0)
    {Supplicant\\{\scriptsize\itshape(Laptop)}};
\node[font=\bfseries\small, align=center, text=textlight] (aut) at (6, 0)
    {Authenticator\\{\scriptsize\itshape(Switch / AP)}};
\node[font=\bfseries\small, align=center, text=textlight] (rad) at (12, 0)
    {RADIUS Server\\{\scriptsize\itshape(FreeRADIUS/ISE)}};

%% Lifeline verticals
\draw[thick, gray!40!white] (0,  -0.4) -- (0,  -13.5);
\draw[thick, gray!40!white] (6,  -0.4) -- (6,  -13.5);
\draw[thick, gray!40!white] (12, -0.4) -- (12, -13.5);

%% Activation boxes
\draw[draw=blue!60!white, fill=blue!15!black] (-0.25,-0.4) rectangle (0.25,-13.5);
\draw[draw=orange!60!white, fill=orange!15!black] (5.75,-0.4) rectangle (6.25,-13.5);
\draw[draw=violet!60!white, fill=violet!20!black] (11.75,-0.4) rectangle (12.25,-13.5);

%% State label: port unauthorized
\node[lbls, align=center, text=orange!70!white] at (6, -1.1) {Puerto:\\UNAUTHORIZED};

%% ── Paso 1: EAPOL-Start ─────────────────────────────────────────────────────
\draw[msg, cyan!70!white] (0.25,-1.6) -- (5.75,-1.6)
    node[midway, above, font=\small, text=cyan!80!white] {\textbf{EAPOL-Start}};
\node[lbls, left] at (0,-1.6) {(1)};

%% ── Paso 2: EAP-Request Identity ────────────────────────────────────────────
\draw[msgr, orange!70!white] (0.25,-2.6) -- (5.75,-2.6)
    node[midway, above, font=\small, text=orange!80!white] {\textbf{EAP-Request / Identity}};
\node[lbls, left] at (0,-2.6) {(2)};

%% ── Paso 3: EAP-Response Identity ───────────────────────────────────────────
\draw[msg, cyan!70!white] (0.25,-3.6) -- (5.75,-3.6)
    node[midway, above, font=\small, text=cyan!80!white] {\textbf{EAP-Response / Identity}};
\node[lbls, left] at (0,-3.6) {(3)};
\node[lbl, below right, font=\scriptsize, text=gray!70!white] at (3,-3.6) {user@empresa.com};

%% ── Paso 4: RADIUS Access-Request ───────────────────────────────────────────
\draw[msg, orange!70!white] (6.25,-4.5) -- (11.75,-4.5)
    node[midway, above, font=\small, text=orange!80!white] {\textbf{RADIUS Access-Request}};
\node[lbls, right] at (12,-4.5) {(4)};
\node[lbl, below right, font=\scriptsize, text=gray!70!white] at (7,-4.5) {EAP-Message encapsulado};

%% ── Paso 5: EAP TLS Tunnel negotiation ──────────────────────────────────────
\draw[msgr, violet!80!white] (6.25,-5.5) -- (11.75,-5.5)
    node[midway, above, font=\small, text=violet!90!white] {\textbf{Access-Challenge (TLS ServerHello)}};
\node[lbls, right] at (12,-5.5) {(5)};

\draw[msgr, orange!70!white] (0.25,-6.3) -- (5.75,-6.3)
    node[midway, above, font=\small, text=orange!80!white] {\textbf{EAP-Request (TLS ServerHello)}};

\draw[msg, cyan!70!white] (0.25,-7.1) -- (5.75,-7.1)
    node[midway, above, font=\small, text=cyan!80!white] {\textbf{EAP-Response (ClientCert + Finished)}};

\draw[msg, orange!70!white] (6.25,-7.9) -- (11.75,-7.9)
    node[midway, above, font=\small, text=orange!80!white] {\textbf{Access-Request (ClientCert)}};

\node[lbls, text=yellow!70!white] at (6,-8.6) {\emoji{locked} Túnel TLS establecido};
\node[lbls] at (6,-9.0) {RADIUS verifica certificado X.509 del cliente};

%% ── Paso 6: RADIUS Access-Accept ────────────────────────────────────────────
\draw[msgr, green!60!white, line width=1.5pt] (6.25,-9.7) -- (11.75,-9.7)
    node[midway, above, font=\small\bfseries, text=green!70!white] {\textbf{RADIUS Access-Accept}};
\node[lbls, right] at (12,-9.7) {(6)};
\node[lbl, below right, font=\scriptsize, text=green!60!white] at (7,-9.7)
    {VLAN=10, Session-Timeout=3600};

%% ── Paso 7: EAP-Success ─────────────────────────────────────────────────────
\draw[msgr, green!60!white, line width=1.5pt] (0.25,-10.6) -- (5.75,-10.6)
    node[midway, above, font=\small\bfseries, text=green!70!white] {\textbf{EAP-Success}};
\node[lbls, left] at (0,-10.6) {(7)};

%% ── Paso 8: Puerto autorizado ───────────────────────────────────────────────
\node[lbls, align=center, text=green!60!white] at (6,-11.4) {Puerto:\\AUTHORIZED};
\node[font=\small, text=green!60!white, align=center] at (6,-12.3)
    {\textbf{Cliente asignado a VLAN 10}\\
     \textbf{Tráfico de datos permitido}};

%% Caption note
\node[font=\scriptsize\itshape, text=gray!60!white, align=center] at (6,-13.3)
    {EAPoL (EAP over LAN) entre Supplicant y Authenticator |
     RADIUS/UDP entre Authenticator y Server};

\end{tikzpicture}
\end{center}

\begin{moderno}
\textbf{Google BeyondCorp --- El dispositivo como identidad}

La implementación de 802.1X en Google va más allá del usuario. En el modelo
\textbf{BeyondCorp} (el origen del término ``Zero Trust''), la autenticación
de red usa \textbf{EAP-TLS donde el certificado pertenece al DISPOSITIVO},
no al usuario:

\begin{itemize}
    \item Cada laptop, teléfono y servidor administrado por Google recibe un
          \textbf{certificado X.509 único} emitido por la PKI interna de Google.
    \item El 802.1X valida ese certificado de dispositivo, no una contraseña de usuario.
    \item Si un empleado trae su laptop personal (\textit{dispositivo no administrado}),
          el AP la rechaza en el 802.1X --- no hay RADIUS Access-Accept ---
          independientemente de que el empleado tenga credenciales válidas de usuario.
    \item \textbf{Credenciales de usuario robadas + dispositivo no administrado
          = cero acceso a la red interna}.
\end{itemize}

\emoji{brain} Este modelo invierte el supuesto de confianza tradicional:
en vez de ``confío en quien está dentro de la red'', el modelo es
``\textit{no confío en nadie hasta verificar dispositivo, usuario y contexto}''.
La red es ahora solo un canal de transporte --- el acceso real lo controla
la capa de aplicación con certificados y políticas de dispositivo.
\end{moderno}

% chapters_parte2/02_switching_avanzado.tex

\chapter{\eh{electric-plug} Switching Avanzado}

Si en la Parte 1 aprendiste que un switch aprende MACs y reenvía tramas, aquí vas a aprender
cómo los switches de nivel empresarial manejan escenarios que un switch hogareño ni soñaría:
múltiples instancias de Spanning Tree, cifrado en el cable, VLANs anidadas, aislamiento de
puertos, control de tormentas de tráfico, y autenticación de dispositivos en el puerto físico.

% ─────────────────────────────────────────────
\section{\eh{repeat-button} MSTP — Multiple Spanning Tree Protocol (802.1s)}

\begin{concepto}
\textbf{Problema con RSTP clásico:} una sola instancia STP para \emph{todas} las VLANs.
Eso significa que si el árbol bloquea un enlace, lo bloquea para \textbf{todo el tráfico},
sin importar la VLAN. Es como cerrar un carril de una autopista y afectar a todos los
vehículos, aunque la mitad fueran por destinos distintos.

\textbf{MSTP (802.1s)} mapea grupos de VLANs a instancias STP independientes, cada una
con su propio árbol y su propio Root Bridge. Resultado: diferentes VLANs pueden usar
diferentes caminos físicos de forma simultánea.
\end{concepto}

\subsection{Conceptos clave}

\begin{itemize}
    \item \textbf{MST Region:} conjunto de switches con la misma configuración MST
          (mismo nombre de región, número de revisión y mapeo VLAN$\to$instancia). Los
          switches fuera de la región se ven como un único switch virtual (boundary).
    \item \textbf{IST (Internal Spanning Tree / Instancia 0):} instancia especial que
          interopera con switches RSTP/STP clásicos en la frontera de la región.
    \item \textbf{MSTI (MST Instance):} instancias 1--64, cada una con su propio Root
          Bridge y topología independiente.
\end{itemize}

\subsection{Configuración Cisco IOS}

\begin{codigo}
\begin{lstlisting}
spanning-tree mode mst

spanning-tree mst configuration
  name CORP_REGION
  revision 1
  instance 1 vlan 10-20
  instance 2 vlan 30-40
  exit

! Root bridge para instancia 1
spanning-tree mst 1 priority 4096

! Root bridge para instancia 2
spanning-tree mst 2 priority 4096
\end{lstlisting}
\end{codigo}

\subsection{Diagrama: MSTP balanceando tráfico}

\begin{center}
\begin{tikzpicture}[
    sw/.style={draw=cyan!60!white, fill=darkcardgray, text=textlight,
               rectangle, rounded corners=4pt, minimum width=2.4cm,
               minimum height=0.8cm, font=\small\bfseries, align=center},
    fwd1/.style={draw=blue!50!white, ultra thick},
    fwd2/.style={draw=green!50!white, ultra thick},
    blk/.style={thick, dashed, draw=orange!60!white},
]
    \node[sw] (a) at (0,3)   {SW-A\\Root MST1};
    \node[sw] (b) at (4,3)   {SW-B\\Root MST2};
    \node[sw] (c) at (2,1.5) {SW-C};
    \node[sw] (d) at (0,0)   {SW-D};
    \node[sw] (e) at (4,0)   {SW-E};

    \draw[fwd1] (a) -- node[left, font=\tiny, text=blue!60!white] {VLAN 10-20} (c);
    \draw[fwd2] (b) -- node[right, font=\tiny, text=green!60!white] {VLAN 30-40} (c);
    \draw[fwd1] (c) -- (d);
    \draw[fwd2] (c) -- (e);
    \draw[blk]  (a) -- node[below left, font=\tiny, text=orange!70!white] {BLK} (d);
    \draw[blk]  (b) -- node[below right, font=\tiny, text=orange!70!white] {BLK} (e);
\end{tikzpicture}

{\small VLAN 10-20 usan el camino izquierdo; VLAN 30-40 el derecho.
Cada instancia bloquea enlaces distintos: \textbf{cero desperdicio de ancho de banda}.}
\end{center}

% ─────────────────────────────────────────────
\section{\eh{locked} MACsec --- MAC Security (802.1AE)}

\begin{concepto}
\textbf{MACsec} cifra el tráfico \emph{hop-by-hop} en Layer 2: entre dos switches
adyacentes, entre un switch y una laptop, entre dos datacenters unidos por fibra oscura.

Analogía: TLS es el sobre sellado que va de tu laptop al servidor de Google.
MACsec es el vehículo blindado que transporta ese sobre entre cada par de edificios.
TLS protege el contenido; MACsec protege el cable físico en cada tramo.
\end{concepto}

\subsection{Cómo funciona}

\begin{enumerate}
    \item \textbf{MKA (MACsec Key Agreement):} protocolo basado en 802.1X que negocia
          claves simétricas entre los peers. Corre sobre el mismo puerto que después cifrará.
    \item \textbf{SecY (Security Entity):} bloque hardware/software que cifra/descifra
          cada frame Ethernet en tiempo real.
    \item \textbf{Overhead:} 32 bytes por frame (16 bytes SecTAG + 16 bytes ICV). Latencia
          adicional: microsegundos en hardware con ASIC dedicado.
\end{enumerate}

\subsection{Formato del frame MACsec}

\begin{verbatim}
+----------+----------+------------+-----------------+----------+----------+
| Dst MAC  | Src MAC  | 0x88E5     |    SecTAG       | Payload  |   ICV    |
| 6 bytes  | 6 bytes  | (S-Tag ET) |  8-16 bytes     | cifrado  | 16 bytes |
+----------+----------+------------+-----------------+----------+----------+
                                    TCI|AN|SL|PN
\end{verbatim}

\begin{cuidado}
MACsec requiere soporte en \textbf{hardware} (ASIC). No todos los switches lo implementan.
Cisco Catalyst 9000 y Nexus 9000 lo soportan nativamente. Switches de gama baja, no.
\end{cuidado}

\subsection{TLS vs MACsec}

\begin{center}
\begin{tabular}{lll}
\toprule
\textbf{Característica} & \textbf{TLS} & \textbf{MACsec} \\
\midrule
Capa           & L7 (aplicación)   & L2 (enlace) \\
Alcance        & End-to-end        & Hop-by-hop \\
Qué cifra      & Payload TCP/UDP   & Frame Ethernet completo \\
Visible al ISP & Solo IPs          & Nada (MACs también ocultas) \\
Requiere app   & Sí                & No (transparente a la app) \\
\bottomrule
\end{tabular}
\end{center}

% ─────────────────────────────────────────────
\section{\eh{nut-and-bolt} QinQ --- Double Tagging (802.1ad)}

\begin{concepto}
\textbf{Problema:} un cliente tiene sus propias VLANs (1--4094). El carrier (ISP) también
usa VLANs internamente. Si ambos usan VLAN 100, hay colisión: el switch del carrier
no distingue ``VLAN 100 del Cliente A'' de ``VLAN 100 del Cliente B''.

\textbf{QinQ} añade un \emph{segundo tag} de VLAN. El tag interno (C-Tag, Customer Tag)
pertenece al cliente. El tag externo (S-Tag, Service Tag) pertenece al carrier.
El carrier solo ve S-Tags; los C-Tags son opacos para él.
\end{concepto}

\subsection{Formato del frame QinQ}

\begin{verbatim}
+---------+---------+----------+---------+----------+---------+
| Dst MAC | Src MAC | 0x88A8   | S-Tag   | 0x8100   | C-Tag   |
| 6 bytes | 6 bytes |(S-Tag ET)|(4 bytes)|(C-Tag ET)|(4 bytes)|
+---------+---------+----------+---------+----------+---------+
          ^-- Carrier procesa hasta aquí
                                          ^-- Cliente ve esto
\end{verbatim}

\begin{ejemplo}
Telmex te vende metro ethernet para conectar Polanco con Vallejo. Internamente usa S-Tag
500 para identificar tu tráfico. Dentro de tu red usas VLAN 10 (datos) y VLAN 20 (voz).
Los switches de Telmex no saben nada de tus VLANs internas: solo mueven frames con S-Tag 500.
\end{ejemplo}

% ─────────────────────────────────────────────
\section{\eh{house} Private VLANs}

\begin{concepto}
\textbf{Private VLANs (PVLANs)} resuelven el aislamiento \emph{dentro} de una VLAN.
Escenario: hosting con 50 VMs del Cliente A en la misma VLAN. No quieres que se vean
entre sí — solo deben alcanzar el gateway. Con VLANs normales, cualquier VM habla con
cualquier otra. PVLANs lo impiden a nivel de switch, sin necesidad de firewall adicional.
\end{concepto}

\subsection{Tipos de puertos}

\begin{center}
\begin{tabular}{lll}
\toprule
\textbf{Tipo} & \textbf{Puede hablar con} & \textbf{Caso de uso} \\
\midrule
Promiscuous  & Todos                        & Router / gateway / firewall \\
Isolated     & Solo Promiscuous             & VM aislada, servidor de cliente \\
Community    & Su comunidad + Promiscuous   & VMs del mismo tenant que sí se comunican \\
\bottomrule
\end{tabular}
\end{center}

\begin{codigo}
\begin{lstlisting}
vlan 100
  private-vlan primary
vlan 101
  private-vlan isolated
vlan 102
  private-vlan community
vlan 100
  private-vlan association 101,102

! Puerto promiscuous (router/gateway)
interface GigabitEthernet0/1
  switchport mode private-vlan promiscuous
  switchport private-vlan mapping 100 101,102

! Puerto isolated (VM aislada)
interface GigabitEthernet0/2
  switchport mode private-vlan host
  switchport private-vlan host-association 100 101
\end{lstlisting}
\end{codigo}

% ─────────────────────────────────────────────
\section{\eh{satellite-antenna} IGMP Snooping}

\begin{concepto}
Sin IGMP Snooping, un switch trata tráfico multicast como \textbf{broadcast}: lo manda
a todos los puertos del segmento. Un stream de video 4K a 20 Mbps llega a todos,
aunque solo una laptop lo haya pedido.

Con IGMP Snooping, el switch \emph{escucha} los mensajes IGMP que cruzan sus puertos
y construye una tabla de ``qué puerto quiere qué grupo multicast''. El tráfico
solo va a los puertos interesados.
\end{concepto}

\subsection{Flujo de aprendizaje}

\begin{enumerate}
    \item Router manda \textbf{IGMP General Query} a \texttt{224.0.0.1} cada 60 s.
    \item Host interesado responde con \textbf{IGMP Membership Report} (grupo específico).
    \item Switch anota: ``Puerto X quiere grupo 239.1.1.1''.
    \item Host que ya no quiere el grupo manda \textbf{IGMP Leave}.
    \item Switch elimina la entrada (o espera timeout de membresía).
\end{enumerate}

\begin{cuidado}
IGMP Snooping necesita un \textbf{IGMP Querier} en la VLAN (normalmente el router
multicast). Sin Querier, los hosts nunca envían Reports espontáneamente y el switch
no aprende → el tráfico multicast termina siendo flooded de todas formas.

Si no hay router multicast, configura \texttt{ip igmp snooping querier} en el switch.
\end{cuidado}

% ─────────────────────────────────────────────
\section{\eh{cloud-with-lightning} Storm Control}

\begin{concepto}
Un loop de red o NIC defectuoso puede generar tráfico broadcast/multicast que se multiplica
exponencialmente: un frame genera 2, esos 2 generan 4, y en segundos tienes una
\emph{broadcast storm} que satura todos los enlaces al 100\%.

\textbf{Storm Control} monitorea el porcentaje de ancho de banda consumido por
broadcast/multicast/unknown unicast. Si supera el umbral, descarta los frames en exceso
(o apaga el puerto).
\end{concepto}

\begin{codigo}
\begin{lstlisting}
interface GigabitEthernet0/1
  storm-control broadcast level 20.00
  storm-control multicast level 10.00
  storm-control action shutdown
\end{lstlisting}
\end{codigo}

\begin{center}
\begin{tabular}{ll}
\toprule
\textbf{Acción} & \textbf{Comportamiento} \\
\midrule
\texttt{drop}     & Descarta exceso, puerto sigue activo \\
\texttt{shutdown} & Apaga el puerto (err-disabled) \\
\texttt{trap}     & Genera SNMP trap al NMS \\
\bottomrule
\end{tabular}
\end{center}

\begin{itemize}
    \item \textbf{STP} previene loops $\to$ evita que la tormenta ocurra.
    \item \textbf{Storm Control} limita el daño si la tormenta \emph{ya ocurrió}.
    Son complementarios, no alternativos.
\end{itemize}

% ─────────────────────────────────────────────
\section{\eh{key} 802.1X en Switches --- Wired NAC}

\begin{concepto}
El mismo 802.1X del Wi-Fi empresarial aplica a \textbf{puertos físicos de switch}. Un
puerto en estado \emph{unauthorized} no permite tráfico de datos hasta que el dispositivo
autentica contra RADIUS.

Analogía: el puerto del switch es el torniquete de un edificio. Sin credenciales válidas
no entras, aunque físicamente hayas insertado el cable.
\end{concepto}

\subsection{MAB --- MAC Authentication Bypass}

\begin{concepto}
Impresoras, teléfonos IP y cámaras no tienen \emph{supplicant} 802.1X. No pueden
iniciar el handshake EAP.

\textbf{MAB} es el fallback: si el dispositivo no responde al EAPOL-Request inicial,
el switch toma su MAC y la envía a RADIUS como si fuera el ``usuario''. RADIUS tiene
una base de datos de MACs autorizadas.

Limitación: MACs se pueden falsificar (\emph{MAC spoofing}). MAB no es autenticación
fuerte; es ``mejor que nada'' para dispositivos sin 802.1X.
\end{concepto}

\begin{codigo}
\begin{lstlisting}
aaa new-model
aaa authentication dot1x default group radius
aaa authorization network default group radius
dot1x system-auth-control

interface GigabitEthernet0/2
  switchport mode access
  switchport access vlan 10
  authentication port-control auto
  authentication order dot1x mab
  authentication priority dot1x mab
  dot1x pae authenticator
  mab
\end{lstlisting}
\end{codigo}

\subsection{VLANs de contingencia}

\begin{itemize}
    \item \textbf{Guest VLAN:} dispositivos que no inician 802.1X caen aquí (acceso limitado).
    \item \textbf{Auth-Fail VLAN:} intentaron autenticar pero fallaron (portal de error).
    \item \textbf{Critical VLAN:} si RADIUS está caído, el puerto cae aquí en vez de
          bloquearse por completo. Evita lockout total de la red.
\end{itemize}

\begin{moderno}
\emoji{sparkles} \textbf{Cisco ISE (Identity Services Engine)} integra 802.1X + MAB +
postura de seguridad (¿el dispositivo tiene antivirus? ¿parches aplicados?) + profiling
automático (identifica si es impresora HP o iPhone) y asigna políticas por identidad,
no por puerto físico. Es el cerebro del NAC empresarial moderno.
\end{moderno}

% chapters_parte2/03_sdn_openflow.tex

\chapter{\eh{control-knobs} SDN y OpenFlow}

Imagina que cada semáforo de una ciudad toma sus propias decisiones de forma independiente:
no sabe qué pasa en el resto de la ciudad, no puede coordinarse con los demás, y para
cambiar una política necesitas ir físicamente a cada semáforo. Eso es el networking
tradicional. \textbf{SDN} pone un controlador central que tiene visibilidad de toda la
red y puede programar cada ``semáforo'' (switch) desde un solo punto.

% ─────────────────────────────────────────────
\section{\eh{warning} El Problema del Networking Tradicional}

\begin{concepto}
En un router o switch tradicional, el \textbf{control plane} (quién decide a dónde va
el paquete) y el \textbf{data plane} (quién realmente lo reenvía) viven \emph{en el
mismo dispositivo}. Esto implica:

\begin{itemize}
    \item Para cambiar política de routing, hay que reconfigurar cada dispositivo
          individualmente (SSH a cada switch, cambiar ACLs, redistribuir rutas).
    \item No hay visibilidad global: cada router solo conoce su vecindad.
    \item Nuevas funciones requieren actualizar firmware en \emph{todos} los dispositivos.
    \item La red no puede adaptarse automáticamente a condiciones cambiantes.
\end{itemize}
\end{concepto}

% ─────────────────────────────────────────────
\section{\eh{building-construction} Arquitectura SDN}

\begin{concepto}
\textbf{SDN (Software-Defined Networking)} separa el control plane del data plane.
El control plane se centraliza en un \textbf{SDN Controller}; los switches solo hacen
forwarding según las instrucciones del controlador.
\end{concepto}

\subsection{Las tres capas}

\begin{center}
\begin{tikzpicture}[
    layer/.style={draw=cyan!50!white, fill=darkcardgray, text=textlight,
                  rectangle, rounded corners=6pt, minimum width=8cm,
                  minimum height=1.2cm, font=\bfseries, align=center},
    api/.style={draw=yellow!60!white, text=yellow!80!white, font=\small\itshape},
]
    \node[layer] (app)  at (0,4.5) {Application Layer\\
        {\small\normalfont Load balancing, firewall, monitoreo, orquestación}};
    \node[layer] (ctrl) at (0,2.5) {Control Layer (SDN Controller)\\
        {\small\normalfont ONOS, OpenDaylight, Ryu}};
    \node[layer] (infra) at (0,0.5) {Infrastructure Layer\\
        {\small\normalfont Switches / Routers (solo data plane)}};

    \draw[-{Stealth}, thick, draw=yellow!60!white]
        (app.south) -- node[api, right] {Northbound API (REST/gRPC)} (ctrl.north);
    \draw[-{Stealth}, thick, draw=orange!60!white]
        (ctrl.south) -- node[api, right] {Southbound API (OpenFlow/NETCONF)} (infra.north);
\end{tikzpicture}
\end{center}

\subsection{Northbound vs Southbound API}

\begin{center}
\begin{tabular}{lll}
\toprule
\textbf{API} & \textbf{Entre} & \textbf{Protocolos típicos} \\
\midrule
Northbound & Aplicaciones $\to$ Controlador & REST, gRPC, GraphQL \\
Southbound & Controlador $\to$ Switches     & OpenFlow, NETCONF, gRPC/gNMI \\
East/West  & Controlador $\to$ Controlador  & OVSDB, clustering protocols \\
\bottomrule
\end{tabular}
\end{center}

% ─────────────────────────────────────────────
\section{\eh{scroll} OpenFlow Protocol}

\begin{concepto}
\textbf{OpenFlow} es el protocolo southbound más extendido. El controlador instala
\emph{flow entries} en la \textbf{flow table} de cada switch OpenFlow. Cuando llega
un paquete, el switch lo compara contra las entries y ejecuta las acciones indicadas.

Analogía: el switch es un mesero con un menú fijo (flow table). El controlador es el
chef que decide qué platos ofrece ese mesero. El mesero no improvisa: si el paquete
no está en el menú, lo manda al chef (controller) para que decida.
\end{concepto}

\subsection{Flow Entry: Match + Actions}

Una flow entry tiene dos partes:

\begin{itemize}
    \item \textbf{Match fields:} los campos del paquete a comparar (puerto de entrada,
          MAC src/dst, VLAN, IP src/dst, protocolo, puerto TCP/UDP, etc.).
    \item \textbf{Actions:} qué hacer si hay match (forward al puerto X, drop, modificar
          campo del paquete, encolar, enviar al controlador).
\end{itemize}

\subsection{Formato de flow entry OpenFlow 1.3}

\begin{lstlisting}
+-----------+----------+---------+---------+----------+----------+
| Match     | Priority | Counters| Timeout | Cookie   | Actions  |
| (12+ campos)         | (pkts,  | (idle,  | (opaco)  | (lista)  |
|           |          |  bytes) |  hard)  |          |          |
+-----------+----------+---------+---------+----------+----------+
\end{lstlisting}

\begin{codigo}
\begin{lstlisting}
# Ejemplo via ovs-ofctl (Open vSwitch)
# Reenviar trafico HTTP (dst port 80) al puerto 2:
ovs-ofctl add-flow br0 \
  "priority=100, ip, nw_proto=6, tp_dst=80, actions=output:2"

# Bloquear IP especifica:
ovs-ofctl add-flow br0 \
  "priority=200, ip, nw_src=10.0.0.5, actions=drop"

# Todo lo demas al controlador:
ovs-ofctl add-flow br0 \
  "priority=0, actions=controller"
\end{lstlisting}
\end{codigo}

\subsection{Pipeline de tablas}

OpenFlow 1.1+ permite múltiples flow tables en secuencia:

\begin{center}
\begin{tikzpicture}[
    tbl/.style={draw=cyan!50!white, fill=darkcardgray, text=textlight,
                rectangle, minimum width=2cm, minimum height=0.9cm,
                font=\small, align=center},
    arr/.style={-{Stealth}, draw=textlight, thick},
]
    \node[tbl] (t0) at (0,0)   {Table 0\\ACL};
    \node[tbl] (t1) at (2.5,0) {Table 1\\VLAN};
    \node[tbl] (t2) at (5,0)   {Table 2\\L3 Route};
    \node[tbl] (t3) at (7.5,0) {Table 3\\QoS};
    \node[tbl, draw=green!50!white] (out) at (10,0) {Output\\Port};

    \draw[arr] (-1.2,0) -- (t0);
    \draw[arr] (t0) -- (t1);
    \draw[arr] (t1) -- (t2);
    \draw[arr] (t2) -- (t3);
    \draw[arr] (t3) -- (out);

    \node[text=yellow!70!white, font=\tiny] at (1.25,-0.8) {GOTO\_TABLE};
    \node[text=yellow!70!white, font=\tiny] at (3.75,-0.8) {GOTO\_TABLE};
    \node[text=yellow!70!white, font=\tiny] at (6.25,-0.8) {GOTO\_TABLE};
    \node[text=yellow!70!white, font=\tiny] at (8.75,-0.8) {GOTO\_TABLE};
\end{tikzpicture}
\end{center}

% ─────────────────────────────────────────────
\section{\eh{desktop-computer} Controladores SDN}

\begin{center}
\begin{tabular}{llll}
\toprule
\textbf{Controlador} & \textbf{Lenguaje} & \textbf{Caso de uso} & \textbf{Nota} \\
\midrule
ONOS        & Java   & Telco, carrier-grade & Clustering nativo \\
OpenDaylight & Java  & Enterprise, modular  & Comunidad grande \\
Ryu         & Python & Testing, académico   & Fácil de extender \\
Floodlight  & Java   & Datacenters          & Fork de NOX \\
\bottomrule
\end{tabular}
\end{center}

\subsection{OVS — Open vSwitch}

\begin{moderno}
\emoji{rocket} \textbf{Open vSwitch (OVS)} es la implementación \emph{software} de un
switch OpenFlow. Vive en el hypervisor (Linux kernel) y conecta VMs entre sí y con la
red física. Es el switch virtual por defecto en OpenStack y muchos entornos cloud privados.

Kubernetes también lo usa: \textbf{OVN (Open Virtual Network)} sobre OVS implementa
el modelo de red de pods con OpenFlow bajo el capó.
\end{moderno}

% ─────────────────────────────────────────────
\section{\eh{brain} Intent-Based Networking (IBN)}

\begin{concepto}
\textbf{IBN} es la evolución de SDN: en vez de programar flows individuales, el admin
declara una \emph{intención de alto nivel} y el sistema la traduce automáticamente a
reglas de bajo nivel.

Ejemplo de intención:
\begin{itemize}
    \item ``Las VMs de producción nunca deben comunicarse con las de desarrollo.''
    \item ``Todo el tráfico de VoIP debe tener latencia máxima de 20 ms.''
    \item ``Los servidores de la VLAN 100 solo acceden a Internet por el firewall A.''
\end{itemize}

El controlador verifica continuamente que la intención se cumpla y la restaura
automáticamente si algo cambia (\emph{closed-loop automation}).
\end{concepto}

\begin{ejemplo}
\textbf{Cisco DNA Center} y \textbf{Cisco Catalyst Center} son las implementaciones
comerciales de IBN de Cisco. \textbf{VMware NSX} hace lo mismo para redes virtualizadas
de datacenters. En ambos casos, el admin define políticas en una GUI y el sistema genera
las ACLs, SGTs (Security Group Tags) y flows de OpenFlow necesarios.
\end{ejemplo}

% ─────────────────────────────────────────────
\section{\eh{memo} SDN en Producción: Realidades}

\begin{cuidado}
\textbf{El controlador es un single point of failure.} Si el controlador cae, los
switches que dependen de él para tomar decisiones pueden dejar de funcionar (o entrar
en modo failopen/failsecure con comportamiento degradado).

Mitigación: clustering de controladores (ONOS tiene clustering nativo), y configurar
los switches para seguir forwarding con las flows cacheadas durante outages breves.
\end{cuidado}

\begin{center}
\begin{tabular}{ll}
\toprule
\textbf{Ventaja SDN} & \textbf{Limitación SDN} \\
\midrule
Visibilidad global de la red    & Controlador = SPOF si no hay HA \\
Cambios de política centralizados & Latencia extra (paquete $\to$ controlador) \\
Automatización programable      & Curva de aprendizaje alta \\
Desacoplamiento HW/SW           & No todos los switches soportan OpenFlow \\
\bottomrule
\end{tabular}
\end{center}

% chapters_parte2/04_datacenter_vxlan_evpn.tex

\chapter{\eh{office-building} Datacenter: VxLAN y EVPN}

Los datacenters modernos tienen un problema que las VLANs clásicas no pueden resolver:
millones de tenants (clientes), VMs que migran entre servidores físicos sin cambiar de IP,
y redes que deben estirarse entre edificios y regiones geográficas. VxLAN + EVPN es la
solución que usan AWS, Google Cloud, Azure y cualquier datacenter serio.

% ─────────────────────────────────────────────
\section{\eh{warning} Limitaciones de VLANs en Datacenter}

\begin{cuidado}
\begin{itemize}
    \item \textbf{4094 VLANs máximo.} Una nube pública tiene millones de clientes, cada
          uno necesita su red aislada. 4094 no alcanza ni para empezar.
    \item \textbf{STP no escala.} Con miles de switches, Spanning Tree converge lentamente
          y bloquea enlaces que podrían estar transmitiendo.
    \item \textbf{VMs migran físicamente.} Cuando vMotion mueve una VM de un servidor a
          otro, necesita mantener la misma IP (mismo L2). Con VLANs, eso requiere que
          ambos servidores estén en el mismo dominio L2, lo que limita la distancia física.
    \item \textbf{L2 sobre L3 es complejo.} Los datacenters rutean internamente (L3 escala),
          pero las aplicaciones esperan L2 (broadcast, ARP). VLANs no resuelven esto.
\end{itemize}
\end{cuidado}

% ─────────────────────────────────────────────
\section{\eh{package} VxLAN --- Virtual Extensible LAN}

\begin{concepto}
\textbf{VxLAN} crea un \emph{overlay} L2 sobre una red L3. Encapsula frames Ethernet
completos dentro de paquetes UDP, permitiendo que dos VMs en servidores físicamente
separados (incluso en datacenters distintos) crean que están en el mismo segmento Ethernet.

\begin{itemize}
    \item \textbf{VNI (VxLAN Network Identifier):} 24 bits $\to$ \textbf{16 millones}
          de segmentos virtuales. Resuelve el límite de 4094 VLANs.
    \item \textbf{VTEP (VxLAN Tunnel Endpoint):} el dispositivo que encapsula/desencapsula.
          Puede ser un switch físico (Leaf) o un servidor Linux con OVS.
    \item \textbf{Puerto UDP:} 4789 (estándar IANA).
\end{itemize}
\end{concepto}

\subsection{Formato del frame VxLAN}

\begin{lstlisting}
+---------------+----------+----------+-----------+------------------+
| Outer Ethernet| Outer IP | Outer UDP| VxLAN Hdr | Inner Ethernet   |
| (VTEP src/dst)|(VTEP IPs)| dst:4789 | (VNI 24b) | Frame original   |
| 14 bytes      | 20 bytes | 8 bytes  | 8 bytes   | variable         |
+---------------+----------+----------+-----------+------------------+
                  ^ Red física underlay              ^ Red virtual overlay
\end{lstlisting}

\begin{itemize}
    \item El frame Ethernet original (inner) viaja intacto dentro del túnel UDP.
    \item La red underlay solo ve IPs de VTEPs; no sabe nada de las VMs.
    \item \textbf{Overhead total:} 50 bytes. En enlaces de 10 GbE, es despreciable.
\end{itemize}

% ─────────────────────────────────────────────
\section{\eh{deciduous-tree} Topología Spine-Leaf}

\begin{concepto}
Las redes de datacenter tradicionales tienen 3 capas: Core $\to$ Distribution $\to$ Access.
Spanning Tree bloquea links redundantes y la latencia varía según el camino.

\textbf{Spine-Leaf} reemplaza esas 3 capas con solo 2, eliminando STP y maximizando
el uso de todos los enlaces mediante ECMP (Equal-Cost Multi-Path).
\end{concepto}

\subsection{Reglas del diseño Spine-Leaf}

\begin{itemize}
    \item Cada \textbf{Leaf} conecta a \textbf{todos} los Spines (ningún Leaf conecta
          directamente a otro Leaf).
    \item Cada \textbf{Spine} conecta a \textbf{todos} los Leaves (ningún Spine conecta
          a otro Spine).
    \item Entre cualquier par de servidores: \textbf{máximo 2 hops} (Leaf $\to$ Spine
          $\to$ Leaf). Latencia predecible y constante.
    \item Para escalar: agregar más Leaves (más puertos de servidor) o más Spines (más
          ancho de banda de east-west).
\end{itemize}

\subsection{Diagrama Spine-Leaf con VxLAN}

\begin{center}
\begin{tikzpicture}[
    spine/.style={draw=cyan!60!white, fill=darkcardviolet, text=textlight,
                  rectangle, rounded corners=4pt, minimum width=2.2cm,
                  minimum height=0.8cm, font=\small\bfseries, align=center},
    leaf/.style={draw=green!50!white, fill=darkcardgreen, text=textlight,
                 rectangle, rounded corners=4pt, minimum width=2.2cm,
                 minimum height=0.8cm, font=\small\bfseries, align=center},
    srv/.style={draw=gray!50!white, fill=darkcardgray, text=textlight,
                rectangle, minimum width=1.5cm, minimum height=0.6cm,
                font=\tiny, align=center},
    link/.style={thick, draw=gray!60!white},
    vtep/.style={text=yellow!70!white, font=\tiny\itshape},
]
    % Spines
    \node[spine] (s1) at (1.5,4) {Spine 1};
    \node[spine] (s2) at (4.5,4) {Spine 2};

    % Leaves
    \node[leaf] (l1) at (0,2) {Leaf 1\\VTEP};
    \node[leaf] (l2) at (3,2) {Leaf 2\\VTEP};
    \node[leaf] (l3) at (6,2) {Leaf 3\\VTEP};

    % Servers
    \node[srv] (srv1) at (-0.7,0.3) {VM-A\\10.0.1.10};
    \node[srv] (srv2) at (0.7,0.3)  {VM-B\\10.0.1.11};
    \node[srv] (srv3) at (2.3,0.3)  {VM-C\\10.0.2.10};
    \node[srv] (srv4) at (3.7,0.3)  {VM-D\\10.0.2.11};
    \node[srv] (srv5) at (5.3,0.3)  {VM-E\\10.0.3.10};
    \node[srv] (srv6) at (6.7,0.3)  {VM-F\\10.0.3.11};

    % Spine-Leaf links
    \draw[link] (l1) -- (s1); \draw[link] (l1) -- (s2);
    \draw[link] (l2) -- (s1); \draw[link] (l2) -- (s2);
    \draw[link] (l3) -- (s1); \draw[link] (l3) -- (s2);

    % Leaf-Server links
    \draw[link] (l1) -- (srv1); \draw[link] (l1) -- (srv2);
    \draw[link] (l2) -- (srv3); \draw[link] (l2) -- (srv4);
    \draw[link] (l3) -- (srv5); \draw[link] (l3) -- (srv6);

    % VxLAN tunnel annotation
    \draw[draw=yellow!60!white, dashed, thick, bend left=20]
        (l1) to node[vtep, above] {VxLAN tunnel VNI=100} (l3);
\end{tikzpicture}

{\small VM-A y VM-E están en el mismo VNI 100 (misma red virtual) aunque en Leaves distintos.
El tráfico viaja encapsulado en UDP: Leaf1 $\to$ Spine $\to$ Leaf3.}
\end{center}

% ─────────────────────────────────────────────
\section{\eh{satellite-antenna} EVPN --- Ethernet VPN}

\begin{concepto}
VxLAN sola tiene un problema: ¿cómo sabe un VTEP a qué otro VTEP enviar un frame si
no conoce la MAC destino? La respuesta naïve es \textbf{flood-and-learn}: inundar el
BUM traffic (Broadcast, Unknown-unicast, Multicast) a todos los VTEPs y aprender MACs
por el tráfico de respuesta. Eso no escala con miles de VTEPs.

\textbf{EVPN} usa \textbf{BGP} como protocolo de control para distribuir información
MAC/IP entre VTEPs. Cada VTEP anuncia las MACs/IPs de sus VMs locales vía BGP, y todos
los demás VTEPs aprenden esa información \emph{antes} de que llegue el primer paquete.
Resultado: no hay flooding innecesario.
\end{concepto}

\subsection{Route Types de EVPN}

\begin{center}
\begin{tabular}{lll}
\toprule
\textbf{Type} & \textbf{Nombre} & \textbf{Propósito} \\
\midrule
Type 1 & Ethernet A-D Route        & Multihoming, failover rápido \\
Type 2 & MAC/IP Advertisement      & Aprender MAC + IP de VMs remotas \\
Type 3 & Inclusive Multicast Route & Registro para recibir BUM traffic \\
Type 4 & Ethernet Segment Route    & Descubrimiento de multihoming \\
Type 5 & IP Prefix Route           & Routing L3 entre VNIs distintos \\
\bottomrule
\end{tabular}
\end{center}

\begin{ejemplo}
\textbf{Flujo con EVPN:}
\begin{enumerate}
    \item VM-A (en Leaf1) envía ARP Request buscando VM-E.
    \item Leaf1 tiene en su BGP EVPN table: VM-E (MAC=aa:bb:...) está en VTEP Leaf3
          (IP=10.1.1.3), VNI=100. Aprendido vía Type 2 Route.
    \item Leaf1 \textbf{no floodea}. Responde el ARP directamente (ARP suppression)
          con la MAC de VM-E desde su tabla local.
    \item VM-A envía el frame. Leaf1 lo encapsula en VxLAN hacia Leaf3.
    \item Leaf3 desencapsula y entrega a VM-E.
\end{enumerate}
Sin EVPN: el ARP hubiera sido floodeado a todos los Leaves. Con EVPN: cero flooding.
\end{ejemplo}

\subsection{Configuración básica (Cisco NX-OS)}

\begin{codigo}
\begin{lstlisting}
! Habilitar VxLAN y EVPN
feature nv overlay
feature bgp
nv overlay evpn

! NVE interface (VTEP)
interface nve1
  no shutdown
  source-interface loopback0
  member vni 10100
    ingress-replication protocol bgp

! BGP con address-family EVPN
router bgp 65001
  address-family l2vpn evpn
    advertise-pip

neighbor 10.1.1.2 remote-as 65001
  address-family l2vpn evpn
    send-community extended
\end{lstlisting}
\end{codigo}

% ─────────────────────────────────────────────
\section{\eh{up-arrow} Routing Asimétrico vs Simétrico en VxLAN}

\begin{center}
\begin{tabular}{lll}
\toprule
\textbf{Modelo} & \textbf{Dónde rutea} & \textbf{Trade-off} \\
\midrule
Asimétrico & VTEP de ingreso rutea, VTEP de egreso bridgea & Simple pero cada Leaf necesita todos los VNIs \\
Simétrico  & Ambos VTEPs rutean, L3 VNI dedicado           & Más escalable, cada Leaf solo necesita sus VNIs \\
\bottomrule
\end{tabular}
\end{center}

\begin{moderno}
\emoji{rocket} \textbf{VxLAN + EVPN en la nube pública:} AWS usa una variante propietaria
de estos principios internamente. Cada Availability Zone tiene su fabric de Leaf/Spine.
El tráfico entre AZs viaja encapsulado (propietario, no OpenFlow). Los Security Groups
de AWS se implementan sobre estos mismos mecanismos de overlay.
\end{moderno}

% chapters_parte2/05_mpls_qos.tex

\chapter{\eh{railway-track} MPLS y QoS}

El Internet público trata todos los paquetes igual: si hay congestión, se pierden
indiscriminadamente. Eso es inaceptable para voz, video en tiempo real, o los SLAs
de un cliente empresarial que paga por un enlace dedicado. \textbf{MPLS} permite
ingeniería de tráfico precisa; \textbf{QoS} garantiza que el tráfico crítico siempre
tenga recursos.

% ─────────────────────────────────────────────
\section{\eh{label} MPLS --- Multiprotocol Label Switching}

\begin{concepto}
En IP normal, cada router examina la dirección destino y hace una búsqueda en su
tabla de routing (\emph{longest prefix match}) para cada paquete. MPLS reemplaza
esa búsqueda repetida por un sistema de \textbf{etiquetas}: el primer router asigna
una etiqueta numérica al paquete y los routers intermedios solo hacen \emph{swap} de
etiquetas en O(1), sin buscar en tablas de routing completas.

\begin{itemize}
    \item \textbf{LSR (Label Switch Router):} router en el núcleo MPLS que solo hace
          swap de etiquetas.
    \item \textbf{LER (Label Edge Router):} router en el borde que asigna (push) o
          elimina (pop) etiquetas.
    \item \textbf{LSP (Label Switched Path):} el camino predeterminado que sigue un
          flujo de tráfico a través de la red MPLS.
    \item \textbf{FEC (Forwarding Equivalence Class):} conjunto de paquetes que reciben
          el mismo tratamiento de forwarding (misma etiqueta, mismo LSP).
\end{itemize}
\end{concepto}

\subsection{Formato de la etiqueta MPLS}

\begin{lstlisting}
 31       20 19    16 15   8 7        0
+----------+---------+------+----------+
|  Label   |   Exp   |  S   |   TTL    |
| (20 bits)|  (3 b)  | (1b) |  (8 b)   |
+----------+---------+------+----------+
\end{lstlisting}

\begin{itemize}
    \item \textbf{Label (20 bits):} identificador del LSP. Hasta $2^{20} = 1{,}048{,}576$ etiquetas.
    \item \textbf{Exp (3 bits):} Traffic Class para QoS (se mapea a DSCP).
    \item \textbf{S (Stack bit):} 1 = última etiqueta del stack (MPLS permite etiquetas anidadas).
    \item \textbf{TTL (8 bits):} igual que el TTL de IP, evita loops.
\end{itemize}

La etiqueta MPLS se inserta entre el header Ethernet (L2) y el header IP (L3),
por eso se llama protocolo de ``capa 2.5''.

% ─────────────────────────────────────────────
\section{\eh{arrows-counterclockwise} Operaciones de etiquetas}

\begin{center}
\begin{tabular}{lll}
\toprule
\textbf{Operación} & \textbf{Quién la hace} & \textbf{Acción} \\
\midrule
\textbf{Push}  & LER de ingreso  & Agrega etiqueta al paquete IP \\
\textbf{Swap}  & LSR en el núcleo & Reemplaza etiqueta por la siguiente del LSP \\
\textbf{Pop}   & LER de egreso (o penúltimo hop) & Elimina etiqueta, entrega IP normal \\
\bottomrule
\end{tabular}
\end{center}

\begin{ejemplo}
\textbf{Paquete de 192.168.1.1 $\to$ 10.0.0.1 a través de red MPLS de ISP:}
\begin{enumerate}
    \item \textbf{LER-A (ingreso):} reconoce FEC ``destino 10.0.0.0/8'', hace \textbf{push}
          etiqueta 200. Paquete: [L2][MPLS:200][IP:10.0.0.1][\ldots]
    \item \textbf{LSR-1:} etiqueta 200 $\to$ tabla dice swap a 350 hacia LSR-2.
          Paquete: [L2][MPLS:350][IP:10.0.0.1][\ldots]
    \item \textbf{LSR-2:} etiqueta 350 $\to$ swap a 120 hacia LER-B.
    \item \textbf{LER-B (egreso):} hace \textbf{pop}, entrega paquete IP puro.
          Forwarding normal basado en IP hacia 10.0.0.1.
\end{enumerate}
Los LSRs intermedios nunca miraron la IP destino. Solo swapearon etiquetas.
\end{ejemplo}

% ─────────────────────────────────────────────
\section{\eh{wrench} MPLS Traffic Engineering (TE)}

\begin{concepto}
Con routing IP normal, el tráfico siempre toma el camino de menor costo (shortest path).
Si hay un enlace congestionado pero otro libre con mayor costo, IP lo ignora. MPLS-TE
permite \textbf{forzar} LSPs por caminos específicos, balancear carga por múltiples
rutas, y garantizar ancho de banda reservado para flujos críticos.

El protocolo de señalización para MPLS-TE es \textbf{RSVP-TE} (Resource Reservation
Protocol - Traffic Engineering), que reserva recursos a lo largo del LSP.
\end{concepto}

\begin{moderno}
\emoji{rocket} \textbf{Segment Routing (SR-MPLS):} la evolución moderna de MPLS-TE.
En lugar de señalización RSVP hop-a-hop, el router de ingreso codifica el camino
completo como un \emph{stack de etiquetas} (segment list). No hay estado de LSP en
los routers intermedios. Google, Facebook y los grandes ISPs migraron de RSVP-TE
a Segment Routing por su simplicidad operacional.
\end{moderno}

% ─────────────────────────────────────────────
\section{\eh{balance-scale} QoS --- Quality of Service}

\begin{concepto}
\textbf{QoS} es el conjunto de mecanismos que garantiza comportamiento predecible
para tráfico crítico en una red compartida. Sin QoS, todos los paquetes compiten
igual por ancho de banda y, en congestión, cualquier flujo puede perder paquetes.
\end{concepto}

\subsection{Los tres problemas que QoS resuelve}

\begin{cuidado}
\begin{itemize}
    \item \textbf{Bandwidth:} un flujo puede consumir todo el ancho de banda disponible
          (Netflix descargando vs. llamada VoIP en el mismo enlace).
    \item \textbf{Delay (latencia):} colas llenas aumentan la latencia. Voz tolera máximo
          150ms one-way; un backup FTP puede tolerar segundos.
    \item \textbf{Jitter (variación de latencia):} paquetes de voz deben llegar con
          intervalos regulares. Jitter alto causa distorsión de audio.
\end{itemize}
\end{cuidado}

\subsection{Modelo DiffServ y DSCP}

\begin{concepto}
\textbf{DiffServ (Differentiated Services)} es el modelo dominante de QoS en Internet
y redes empresariales. Cada paquete IP lleva un campo \textbf{DSCP (Differentiated
Services Code Point)} de 6 bits en el header IP (parte del campo ToS/Traffic Class).
Los routers aplican tratamiento diferenciado según ese valor, sin mantener estado
por flujo.
\end{concepto}

\begin{center}
\begin{tabular}{llll}
\toprule
\textbf{Clase} & \textbf{DSCP} & \textbf{Valor} & \textbf{Uso típico} \\
\midrule
EF (Expedited Forwarding)   & EF   & 46 (101110) & VoIP, videoconferencia \\
AF41 (Assured Forwarding)   & AF41 & 34 (100010) & Video interactivo \\
AF31                         & AF31 & 26 (011010) & Video streaming \\
AF21                         & AF21 & 18 (010010) & Transacciones críticas \\
AF11                         & AF11 & 10 (001010) & Datos de negocios \\
CS0 (Best Effort)            & BE   &  0 (000000) & Tráfico normal \\
\bottomrule
\end{tabular}
\end{center}

\subsection{Mecanismos de QoS}

\begin{center}
\begin{tikzpicture}[
    box/.style={draw=cyan!50!white, fill=darkcardgray, text=textlight,
                rectangle, rounded corners=4pt, minimum width=3cm,
                minimum height=1cm, font=\small\bfseries, align=center},
    arrow/.style={->, thick, draw=gray!60!white},
]
    \node[box] (class) at (0,0)    {Clasificación\\{\tiny DSCP, ACL, NBAR}};
    \node[box] (mark)  at (4,0)    {Marcado\\{\tiny Set DSCP}};
    \node[box] (queue) at (8,0)    {Colas\\{\tiny LLQ, WFQ, CBWFQ}};
    \node[box] (sched) at (8,-2)   {Scheduling\\{\tiny FIFO, WRR, SP}};
    \node[box] (shape) at (4,-2)   {Shaping/Policing\\{\tiny Token bucket}};
    \node[box] (drop)  at (0,-2)   {Drop Policy\\{\tiny WRED, tail drop}};

    \draw[arrow] (class) -- (mark);
    \draw[arrow] (mark)  -- (queue);
    \draw[arrow] (queue) -- (sched);
    \draw[arrow] (sched) -- (shape);
    \draw[arrow] (shape) -- (drop);
\end{tikzpicture}
\end{center}

\begin{itemize}
    \item \textbf{LLQ (Low Latency Queue):} cola de alta prioridad para VoIP/video.
          Se sirve primero, siempre. Garantiza latencia mínima.
    \item \textbf{CBWFQ (Class-Based WFQ):} asigna porcentajes de ancho de banda
          garantizado por clase. Si nadie más usa su porción, puede tomar el exceso.
    \item \textbf{WRED (Weighted Random Early Detection):} descarta paquetes de baja
          prioridad antes de que la cola se llene, evitando TCP global synchronization.
    \item \textbf{Token Bucket:} modelo para policing/shaping. Se ``acumulan fichas''
          a una tasa CIR (Committed Information Rate); cada paquete consume fichas.
          Sin fichas $\to$ paquete se descarta (policing) o se retrasa (shaping).
\end{itemize}

\begin{ejemplo}
\textbf{QoS en una empresa con 100 Mbps hacia Internet:}
\begin{itemize}
    \item VoIP (DSCP EF): LLQ, 10 Mbps garantizados, latencia $<$10ms.
    \item Videoconferencia (DSCP AF41): CBWFQ, 30 Mbps garantizados.
    \item CRM/ERP (DSCP AF21): CBWFQ, 25 Mbps garantizados.
    \item Navegación general (DSCP BE): 35 Mbps restantes, WRED activo.
\end{itemize}
Si las llamadas VoIP no usan sus 10 Mbps, el resto puede aprovecharlos.
Pero si VoIP necesita esos 10 Mbps, los tiene garantizados sin importar lo demás.
\end{ejemplo}

\begin{moderno}
\emoji{rocket} \textbf{QoS en la nube:} AWS, GCP y Azure ofrecen Enhanced Networking
(SR-IOV) para reducir latencia en instancias. Google Cloud usa \textbf{Andromeda},
su stack de virtualización de red que aplica QoS a nivel de hipervisor. En Kubernetes,
los \texttt{LimitRange} y \texttt{ResourceQuota} son la capa de QoS de aplicación;
el QoS de red real requiere CNI plugins como Cilium con bandwidth enforcement via eBPF.
\end{moderno}

% chapters_parte2/06_kubernetes_networking.tex

\chapter{\eh{gear} Networking en Kubernetes}

Kubernetes es hoy la plataforma estándar de orquestación de contenedores. Su modelo
de red es intencionalmente simple en la interfaz pero complejo en la implementación:
\emph{todo Pod puede hablar con todo Pod sin NAT}. Entender cómo se implementa eso
es fundamental para operar clusters en producción.

% ─────────────────────────────────────────────
\section{\eh{memo} El Modelo de Red de Kubernetes}

\begin{concepto}
Kubernetes impone cuatro reglas fundamentales al networking:

\begin{enumerate}
    \item Cada Pod tiene su propia dirección IP única en el cluster.
    \item Todos los Pods pueden comunicarse entre sí sin NAT.
    \item Los nodos pueden comunicarse con todos los Pods sin NAT.
    \item La IP que ve un Pod de sí mismo es la misma que ven los demás Pods.
\end{enumerate}

Estas reglas simplifican el modelo: las aplicaciones no necesitan saber si están en el
mismo nodo o en nodos distintos. El networking subyacente lo resuelve.
\end{concepto}

\begin{cuidado}
Kubernetes \textbf{define} el modelo pero \textbf{no lo implementa}. La implementación
la provee el \textbf{CNI plugin} (Container Network Interface): Cilium, Calico, Flannel,
Weave, etc. Distintos CNIs tienen distintas capacidades (Network Policies, encryption,
eBPF, multicast, etc.).
\end{cuidado}

% ─────────────────────────────────────────────
\section{\eh{electric-plug} Tipos de red en un cluster}

\begin{center}
\begin{tikzpicture}[
    netbox/.style={draw=cyan!50!white, fill=darkcardgray, text=textlight,
                   rectangle, rounded corners=4pt, minimum width=5cm,
                   minimum height=0.9cm, font=\small\bfseries, align=center},
    node/.style={draw=green!50!white, fill=darkcardgreen, text=textlight,
                 rectangle, rounded corners=4pt, minimum width=3cm,
                 minimum height=0.8cm, font=\small, align=center},
    pod/.style={draw=gray!50!white, fill=darkcardgray, text=textlight,
                rectangle, minimum width=1.5cm, minimum height=0.6cm,
                font=\tiny, align=center},
    arrow/.style={->, thick, draw=gray!60!white},
]
    \node[netbox] (pod_net) at (0, 4.5) {Pod Network (ej. 10.244.0.0/16)};
    \node[netbox] (svc_net) at (0, 3)   {Service Network (ej. 10.96.0.0/12)};
    \node[netbox] (node_net) at (0, 1.5) {Node Network (ej. 192.168.1.0/24)};

    \node[node]  (n1) at (-3, 0) {Node 1\\192.168.1.10};
    \node[node]  (n2) at (3, 0)  {Node 2\\192.168.1.11};
    \node[pod]   (p1) at (-4.2, -1.5) {Pod A\\10.244.0.2};
    \node[pod]   (p2) at (-2.2, -1.5) {Pod B\\10.244.0.3};
    \node[pod]   (p3) at (1.8, -1.5)  {Pod C\\10.244.1.2};
    \node[pod]   (p4) at (3.8, -1.5)  {Pod D\\10.244.1.3};

    \draw[arrow] (n1) -- (p1); \draw[arrow] (n1) -- (p2);
    \draw[arrow] (n2) -- (p3); \draw[arrow] (n2) -- (p4);
    \draw[arrow, dashed, draw=yellow!60!white] (p1) to[bend left=30]
        node[font=\tiny, text=yellow!70!white, above] {CNI overlay} (p3);
\end{tikzpicture}
\end{center}

\begin{itemize}
    \item \textbf{Pod Network:} rango de IPs asignadas a Pods. CNI lo administra.
    \item \textbf{Service Network:} IPs virtuales de Services (ClusterIP). Solo existen
          en iptables/eBPF; ningún dispositivo real las tiene.
    \item \textbf{Node Network:} red física/real de los nodos. Preexiste a Kubernetes.
\end{itemize}

% ─────────────────────────────────────────────
\section{\eh{gear} Services y kube-proxy}

\begin{concepto}
Los Pods son efímeros: mueren y resurgen con nuevas IPs. Un \textbf{Service} es una
abstracción estable: tiene una IP virtual fija (ClusterIP) y un DNS name, y hace
load balancing hacia los Pods seleccionados por su \texttt{selector}.

\textbf{kube-proxy} implementa los Services programando reglas en el kernel del nodo:
\end{concepto}

\begin{center}
\begin{tabular}{lll}
\toprule
\textbf{Modo} & \textbf{Mecanismo} & \textbf{Rendimiento} \\
\midrule
userspace & Proxy en espacio de usuario & Lento, legado \\
iptables  & Reglas NAT en iptables       & Bueno, O(n) rules \\
IPVS      & IP Virtual Server del kernel & Mejor, O(1) lookup \\
eBPF      & Cilium sin kube-proxy        & Óptimo, bypass kernel stack \\
\bottomrule
\end{tabular}
\end{center}

\subsection{Tipos de Service}

\begin{itemize}
    \item \textbf{ClusterIP (default):} IP virtual solo accesible dentro del cluster.
          Uso: comunicación interna entre microservicios.
    \item \textbf{NodePort:} expone el Service en un puerto fijo de cada nodo
          (30000--32767). Accesible desde fuera del cluster via \texttt{NodeIP:NodePort}.
    \item \textbf{LoadBalancer:} solicita un load balancer externo al cloud provider
          (AWS ELB, GCP LB). Asigna IP pública automáticamente.
    \item \textbf{ExternalName:} alias DNS a un FQDN externo. Útil para integrar
          servicios externos como si fueran internos.
    \item \textbf{Headless (ClusterIP: None):} no hay IP virtual; DNS devuelve
          directamente las IPs de los Pods. Usado por StatefulSets (Kafka, Cassandra).
\end{itemize}

% ─────────────────────────────────────────────
\section{\eh{shield} Network Policies}

\begin{concepto}
Por defecto, todos los Pods en un cluster pueden comunicarse entre sí (modelo flat).
\textbf{NetworkPolicy} permite definir reglas de firewall a nivel de Pod: qué pods
pueden conectarse a qué otros pods, en qué puertos y protocolos.

Las NetworkPolicies son aditivas: si no hay política, se permite todo. Si hay al menos
una política que selecciona un Pod, ese Pod solo acepta lo que las políticas permiten.
\end{concepto}

\begin{codigo}
\begin{lstlisting}
# Solo permitir trafico al pod "backend" desde pods con label "tier: frontend"
apiVersion: networking.k8s.io/v1
kind: NetworkPolicy
metadata:
  name: backend-ingress
spec:
  podSelector:
    matchLabels:
      app: backend
  policyTypes:
  - Ingress
  ingress:
  - from:
    - podSelector:
        matchLabels:
          tier: frontend
    ports:
    - protocol: TCP
      port: 8080
\end{lstlisting}
\end{codigo}

\begin{cuidado}
NetworkPolicies requieren que el CNI plugin las soporte. Flannel no las implementa.
Calico, Cilium, y Weave sí. Sin un CNI compatible, los objetos NetworkPolicy se crean
pero no tienen efecto alguno --- la red sigue siendo completamente abierta.
\end{cuidado}

% ─────────────────────────────────────────────
\section{\eh{globe-with-meridians} Ingress y Gateway API}

\begin{concepto}
Los Services de tipo LoadBalancer son costosos (un LB externo por Service). \textbf{Ingress}
consolida el tráfico HTTP/HTTPS de múltiples Services detrás de un solo load balancer,
enrutando por host y path.
\end{concepto}

\begin{codigo}
\begin{lstlisting}
apiVersion: networking.k8s.io/v1
kind: Ingress
metadata:
  name: mi-app-ingress
spec:
  rules:
  - host: api.miempresa.com
    http:
      paths:
      - path: /v1
        pathType: Prefix
        backend:
          service:
            name: api-v1
            port:
              number: 80
      - path: /v2
        pathType: Prefix
        backend:
          service:
            name: api-v2
            port:
              number: 80
\end{lstlisting}
\end{codigo}

\begin{moderno}
\emoji{rocket} \textbf{Gateway API} (sucesor de Ingress): API más expresiva con soporte
nativo para TCP, UDP, gRPC, y modelos multi-tenant. Separa responsabilidades entre el
equipo de infraestructura (GatewayClass, Gateway) y el equipo de aplicación (HTTPRoute,
TCPRoute). Cilium, Istio y Nginx ya la implementan. Ingress quedará deprecated eventualmente.
\end{moderno}

% ─────────────────────────────────────────────
\section{\eh{satellite-antenna} CNI Plugins: Cilium vs Calico}

\begin{center}
\begin{tabular}{lll}
\toprule
\textbf{Feature} & \textbf{Cilium} & \textbf{Calico} \\
\midrule
Data plane         & eBPF                & iptables / eBPF (opcional) \\
Rendimiento        & Excelente           & Bueno \\
Network Policies   & L3/L4/L7 (HTTP, gRPC) & L3/L4 \\
Observabilidad     & Hubble (flujos)     & Limitada \\
Modo sin kube-proxy & Sí (eBPF)          & Parcial \\
BGP support        & Sí                  & Sí (más maduro) \\
Encryp. transparente & WireGuard/IPsec   & WireGuard/IPsec \\
\bottomrule
\end{tabular}
\end{center}

\begin{ejemplo}
\textbf{Por qué Cilium usa eBPF:}
Con iptables, cada Service agrega reglas que el kernel evalúa secuencialmente.
Con 10,000 Services, kube-proxy crea cientos de miles de reglas; la actualización
tarda minutos y la evaluación es O(n). eBPF usa mapas hash: lookup O(1), actualizaciones
atómicas en microsegundos. En clusters grandes ($>$500 nodos), Cilium reduce la latencia
de Service lookup de milisegundos a microsegundos.
\end{ejemplo}

% chapters_parte2/07_cloud_networking.tex

\chapter{\eh{cloud} Cloud Networking}

Las nubes públicas no son simplemente servidores remotos con Internet rápido. Son
redes enteras diseñadas desde cero para escala masiva, multi-tenancy y automatización
via API. Entender su arquitectura diferencia a un ingeniero que \emph{usa} la nube
de uno que \emph{diseña} sobre ella.

% ─────────────────────────────────────────────
\section{\eh{building-construction} VPC --- Virtual Private Cloud}

\begin{concepto}
Una \textbf{VPC (Virtual Private Cloud)} es una red virtual aislada dentro de la
nube pública. Cada cliente tiene su propia VPC con su propio espacio de direcciones IP,
tablas de routing, y políticas de seguridad. El aislamiento entre VPCs lo garantiza
el hypervisor/SDN del proveedor, no hardware dedicado.

Componentes fundamentales (nomenclatura AWS, análogos en GCP/Azure):
\begin{itemize}
    \item \textbf{CIDR block:} rango IP de la VPC (ej. 10.0.0.0/16).
    \item \textbf{Subnets:} subdivisiones de la VPC, asociadas a una zona de disponibilidad.
    \item \textbf{Route Tables:} controlan a dónde va el tráfico de cada subnet.
    \item \textbf{Internet Gateway (IGW):} permite que subnets públicas accedan a Internet.
    \item \textbf{NAT Gateway:} permite que subnets privadas salgan a Internet sin
          ser directamente accesibles.
    \item \textbf{Security Groups:} firewall stateful por instancia (como iptables).
    \item \textbf{NACLs (Network ACLs):} firewall stateless por subnet.
\end{itemize}
\end{concepto}

\subsection{Arquitectura típica multi-tier}

\begin{center}
\begin{tikzpicture}[
    tier/.style={draw=cyan!50!white, fill=darkcardgray, text=textlight,
                 rectangle, rounded corners=4pt, minimum width=7cm,
                 minimum height=0.9cm, font=\small\bfseries, align=center},
    gw/.style={draw=yellow!60!white, fill=darkcardviolet, text=textlight,
               rectangle, rounded corners=4pt, minimum width=2cm,
               minimum height=0.7cm, font=\tiny, align=center},
    arrow/.style={->, thick, draw=gray!60!white},
]
    \node[gw] (igw)  at (0, 6)   {Internet Gateway};
    \node[tier] (pub)  at (0, 4.5) {Subnet Pública: Load Balancer, Bastion};
    \node[gw] (nat)  at (4, 3)   {NAT GW};
    \node[tier] (app)  at (0, 3)   {Subnet Privada: App Servers, APIs};
    \node[tier] (db)   at (0, 1.5) {Subnet Privada: Bases de Datos, Caches};
    \node[tier] (vpc)  at (0, 0)   {VPC: 10.0.0.0/16};

    \draw[arrow] (igw) -- (pub);
    \draw[arrow] (pub) -- (app);
    \draw[arrow] (app) -- (db);
    \draw[arrow] (app) -- (nat);
    \draw[arrow, dashed] (nat) -- node[font=\tiny, text=yellow!70!white, right] {Internet} (igw);
\end{tikzpicture}
\end{center}

\begin{itemize}
    \item Subnet pública: accesible desde Internet (IGW en route table).
    \item Subnet privada de app: no accesible directamente, sale a Internet vía NAT.
    \item Subnet privada de DB: sin acceso a Internet en ninguna dirección.
\end{itemize}

% ─────────────────────────────────────────────
\section{\eh{link} Conectividad entre VPCs y On-Premise}

\subsection{VPC Peering}

\begin{concepto}
Conexión privada entre dos VPCs (mismo o distinto account, misma o distinta región).
El tráfico viaja por la backbone del proveedor, no por Internet. \textbf{No transitivo:}
si A está peered con B, y B con C, A \textbf{no} puede hablar con C a través de B.
\end{concepto}

\subsection{Transit Gateway / Cloud Router}

\begin{concepto}
Solución hub-and-spoke para conectar muchas VPCs y redes on-premise. En lugar de
$N^2$ peerings, cada VPC conecta al Transit Gateway (AWS) o Cloud Router (GCP).
Sí es transitivo. Permite routing más granular y segmentación de tráfico.
\end{concepto}

\subsection{VPN y Direct Connect / Dedicated Interconnect}

\begin{center}
\begin{tabular}{lll}
\toprule
\textbf{Opción} & \textbf{Características} & \textbf{Cuándo usar} \\
\midrule
VPN (IPsec)       & Sobre Internet, encriptado, rápido de desplegar & Dev/test, fallback \\
Direct Connect    & Enlace físico dedicado, bajo jitter, SLA fuerte & Producción crítica \\
PrivateLink       & Acceso a servicios del proveedor sin Internet     & S3, RDS, APIs internas \\
\bottomrule
\end{tabular}
\end{center}

% ─────────────────────────────────────────────
\section{\eh{bar-chart} Load Balancing en la nube}

\begin{concepto}
Los proveedores cloud ofrecen load balancers gestionados que escalan automáticamente.
La elección depende del tipo de tráfico:
\end{concepto}

\begin{center}
\begin{tabular}{llll}
\toprule
\textbf{Tipo} & \textbf{Capa OSI} & \textbf{AWS} & \textbf{GCP} \\
\midrule
L4 TCP/UDP    & Transporte & Network LB (NLB)    & TCP/UDP LB \\
L7 HTTP/HTTPS & Aplicación & Application LB (ALB) & HTTP(S) LB \\
Global anycast & L7 global & CloudFront + ALB     & Global HTTP LB \\
\bottomrule
\end{tabular}
\end{center}

\begin{ejemplo}
\textbf{ALB (Application Load Balancer)} puede enrutar por:
\begin{itemize}
    \item Host header: \texttt{api.ejemplo.com} $\to$ Target Group A;
          \texttt{admin.ejemplo.com} $\to$ Target Group B.
    \item Path: \texttt{/api/v1} $\to$ servicio legacy; \texttt{/api/v2} $\to$ servicio nuevo.
    \item Header HTTP: \texttt{X-Version: beta} $\to$ 10\% canary deployment.
    \item Peso: 90\% tráfico a stable, 10\% a canary (weighted target groups).
\end{itemize}
\end{ejemplo}

% ─────────────────────────────────────────────
\section{\eh{lock} Seguridad de Red en la Nube}

\subsection{Security Groups vs NACLs}

\begin{center}
\begin{tabular}{lll}
\toprule
\textbf{Característica} & \textbf{Security Groups} & \textbf{NACLs} \\
\midrule
Estado         & Stateful (respuesta automática) & Stateless (reglas explícitas ida y vuelta) \\
Nivel          & Instancia (ENI)                 & Subnet \\
Reglas         & Solo allow                      & Allow y deny \\
Evaluación     & Todas las reglas juntas          & Numeradas, primera que aplica \\
Uso típico     & Firewall de host                & Bloqueo de IPs maliciosas \\
\bottomrule
\end{tabular}
\end{center}

\begin{cuidado}
NACLs son stateless: si permites tráfico entrante al puerto 80, necesitas una regla
de salida para los puertos efímeros (1024--65535) para que las respuestas puedan salir.
Con Security Groups no: al ser stateful, la respuesta se permite automáticamente.
\end{cuidado}

\subsection{AWS PrivateLink / GCP Private Service Connect}

\begin{moderno}
\emoji{rocket} Acceso a servicios (S3, DynamoDB, APIs de terceros) sin que el tráfico
salga a Internet. El servicio se expone como un endpoint en tu VPC con IP privada.
El tráfico nunca cruza un IGW. Mitiga riesgo de exfiltración y elimina cargos de
NAT Gateway para tráfico a servicios del proveedor (S3 puede ser millones al mes).
\end{moderno}

% ─────────────────────────────────────────────
\section{\eh{satellite} CDN y Edge Networking}

\begin{concepto}
Una \textbf{CDN (Content Delivery Network)} distribuye contenido estático (imágenes,
JS, CSS, videos) a \textbf{PoPs (Points of Presence)} cercanos a los usuarios finales.
El objetivo: reducir latencia sirviendo desde el PoP más próximo al usuario, no desde
el origen.
\end{concepto}

\begin{center}
\begin{tabular}{lll}
\toprule
\textbf{Provider} & \textbf{CDN} & \textbf{Diferenciador} \\
\midrule
AWS         & CloudFront     & Integración nativa con S3, ALB, Lambda@Edge \\
GCP         & Cloud CDN      & Backbone privado Google, anycast global \\
Cloudflare  & Cloudflare CDN & DDoS protection, Workers edge compute \\
Akamai      & Akamai CDN     & Mayor red de PoPs del mundo \\
\bottomrule
\end{tabular}
\end{center}

\begin{ejemplo}
\textbf{Latencia con y sin CDN para usuario en CDMX accediendo a servidor en Virginia:}
\begin{itemize}
    \item Sin CDN: CDMX $\to$ Internet $\to$ Virginia. RTT $\approx$ 80ms. Carga HTML+JS+CSS
          requiere múltiples round-trips: $>$500ms.
    \item Con CDN (PoP en CDMX): JS/CSS servidos localmente. RTT al PoP $\approx$ 5ms.
          Solo llamadas a API van a Virginia. Carga percibida: $<$100ms.
\end{itemize}
\end{ejemplo}

% chapters_parte2/08_multicast.tex

\chapter{\eh{mega} Multicast}

Imagina transmitir un partido de fútbol a 1 millón de espectadores. Con \textbf{unicast},
el servidor enviaría 1 millón de copias del mismo stream: imposible. Con
\textbf{broadcast}, llegaría a todos los dispositivos de la red, incluso a quienes no
quieren verlo. \textbf{Multicast} resuelve esto elegantemente: el servidor envía
\emph{una sola copia} que la red replica solo hacia quienes la solicitaron.

% ─────────────────────────────────────────────
\section{\eh{bar-chart} Unicast vs Broadcast vs Multicast}

\begin{center}
\begin{tabular}{llll}
\toprule
\textbf{Modelo} & \textbf{Destinatario} & \textbf{Copias en red} & \textbf{Uso} \\
\midrule
Unicast   & Un host específico       & Una por destino         & HTTP, SSH, DNS \\
Broadcast & Todos en el segmento L2  & Una (pero llega a todos) & ARP, DHCP \\
Multicast & Grupo de interesados     & Una (la red la replica) & IPTV, routing protocols \\
Anycast   & El más cercano del grupo & Una                     & DNS raíz, CDN \\
\bottomrule
\end{tabular}
\end{center}

\begin{cuidado}
Multicast tiene reputación de ser complejo de implementar y operar. No es mito:
requiere soporte en cada router del camino, gestión de grupos, y protocolos adicionales
(PIM, IGMP, MSDP). Por eso en la nube pública casi no existe multicast nativo: AWS
y GCP no lo soportan en VPCs. Se reemplaza con overlay multicast (software) o CDN.
\end{cuidado}

% ─────────────────────────────────────────────
\section{\eh{label} Direcciones Multicast}

\begin{concepto}
IPv4 reserva el rango \textbf{224.0.0.0/4} para multicast ($224.x.x.x$ a $239.x.x.x$).
Cada \emph{grupo multicast} tiene una dirección en este rango. Un host se ``une'' al
grupo para recibir tráfico; puede abandonarlo en cualquier momento.
\end{concepto}

\begin{center}
\begin{tabular}{lll}
\toprule
\textbf{Rango} & \textbf{Nombre} & \textbf{Uso} \\
\midrule
224.0.0.0/24   & Link-local       & Protocolos de routing (OSPF: 224.0.0.5/6, HSRP: 224.0.0.2) \\
224.0.1.0/24   & Internetwork     & NTP: 224.0.1.1, uso general \\
232.0.0.0/8    & Source-Specific  & SSM (Source-Specific Multicast) \\
239.0.0.0/8    & Admin-scoped     & Multicast privado dentro de organización \\
\bottomrule
\end{tabular}
\end{center}

\subsection{MAC Multicast}

\begin{concepto}
Cuando un router L2 recibe un paquete multicast IP, lo encapsula en un frame Ethernet
con una MAC de destino especial. El bit I/G del primer byte de la MAC indica multicast.
Para IPs multicast, la MAC se construye así:

\begin{center}
\texttt{01:00:5E:[0][últimos 23 bits de la IP multicast]}
\end{center}

Ej: IP \texttt{224.1.1.1} $\to$ MAC \texttt{01:00:5E:01:01:01}.

Consecuencia: los switches L2 no hacen multicast por defecto --- inundan a todos los
puertos como si fuera broadcast. \textbf{IGMP Snooping} soluciona esto.
\end{concepto}

% ─────────────────────────────────────────────
\section{\eh{ear-of-corn} IGMP --- Gestión de Grupos}

\begin{concepto}
\textbf{IGMP (Internet Group Management Protocol)} es el protocolo entre hosts y el
router de primer salto para gestionar membresía a grupos multicast.
\end{concepto}

\begin{center}
\begin{tabular}{lll}
\toprule
\textbf{Versión} & \textbf{Característica clave} & \textbf{Uso} \\
\midrule
IGMPv1 & Join básico, no hay Leave & Obsoleto \\
IGMPv2 & Leave Group message      & Aún común \\
IGMPv3 & Source filtering (SSM)   & Moderno, requerido para SSM \\
\bottomrule
\end{tabular}
\end{center}

\begin{itemize}
    \item \textbf{IGMP Query:} el router pregunta ``¿alguien quiere grupos multicast?''
          cada 60--125 segundos.
    \item \textbf{IGMP Report:} el host responde con los grupos a los que pertenece.
    \item \textbf{IGMP Leave:} el host notifica que abandona un grupo (IGMPv2+).
    \item \textbf{IGMP Snooping:} los switches L2 ``escuchan'' mensajes IGMP y aprenden
          qué puerto tiene interesados en cada grupo, evitando flooding innecesario.
\end{itemize}

% ─────────────────────────────────────────────
\section{\eh{deciduous-tree} PIM --- Routing Multicast entre Routers}

\begin{concepto}
IGMP maneja el último salto (host $\leftrightarrow$ router). Para que el tráfico
multicast fluya entre routers en la red, se usa \textbf{PIM (Protocol Independent
Multicast)}. ``Protocol Independent'' porque funciona sobre cualquier protocolo de
routing unicast (OSPF, BGP, etc.), usando su tabla de routing para construir árboles
de distribución multicast.
\end{concepto}

\subsection{PIM Dense Mode (PIM-DM)}

\begin{itemize}
    \item Asume que todos los segmentos quieren el tráfico multicast.
    \item \textbf{Flood and prune:} inunda todo, luego los routers sin interesados
          envían mensajes Prune para cortar el árbol.
    \item Simple pero ineficiente. Solo adecuado para redes pequeñas donde la mayoría
          de los segmentos sí tienen receptores.
\end{itemize}

\subsection{PIM Sparse Mode (PIM-SM)}

\begin{concepto}
Modelo opuesto: asume que \emph{pocos} segmentos quieren el tráfico. Los routers
deben suscribirse explícitamente. Usa un \textbf{RP (Rendezvous Point)}: router central
donde fuentes y receptores se ``encuentran''.
\end{concepto}

\begin{center}
\begin{tikzpicture}[
    rtr/.style={draw=cyan!50!white, fill=darkcardgray, text=textlight,
                circle, minimum size=0.9cm, font=\tiny\bfseries},
    rp/.style={draw=yellow!60!white, fill=darkcardviolet, text=textlight,
               circle, minimum size=0.9cm, font=\tiny\bfseries},
    host/.style={draw=gray!50!white, fill=darkcardgray, text=textlight,
                 rectangle, minimum width=1.2cm, minimum height=0.6cm, font=\tiny, align=center},
    arrow/.style={->, thick, draw=gray!60!white},
    joinpath/.style={->, thick, draw=green!60!white, dashed},
    srcpath/.style={->, thick, draw=yellow!60!white},
]
    \node[rp]  (rp) at (3, 4)   {RP};
    \node[rtr] (r1) at (0, 2)   {R1};
    \node[rtr] (r2) at (3, 2)   {R2};
    \node[rtr] (r3) at (6, 2)   {R3};
    \node[host] (src) at (0, 0.5) {Fuente\\224.1.1.1};
    \node[host] (rx1) at (3, 0.5) {Receptor 1};
    \node[host] (rx2) at (6, 0.5) {Receptor 2};
    \draw[arrow] (r1) -- (rp);
    \draw[arrow] (r2) -- (rp);
    \draw[arrow] (r3) -- (rp);
    \draw[srcpath] (src) -- node[font=\tiny, text=yellow!70!white, left] {stream} (r1);
    \draw[joinpath] (rx1) -- node[font=\tiny, text=green!70!white, right] {Join} (r2);
    \draw[joinpath] (rx2) -- node[font=\tiny, text=green!70!white, right] {Join} (r3);
    \draw[srcpath] (rp) -- (r2);
    \draw[srcpath] (rp) -- (r3);
\end{tikzpicture}

{\small Receptores hacen Join hacia el RP. Fuente registra su stream en el RP.
RP distribuye a los routers con receptores. Después puede formarse un Shortest Path Tree
directo fuente $\to$ receptor (SPT switchover).}
\end{center}

\subsection{PIM Source-Specific Multicast (PIM-SSM)}

\begin{concepto}
Variante moderna que elimina el RP. Los receptores se suscriben a un grupo multicast
\emph{y} una fuente específica: \texttt{(S, G)} en lugar de \texttt{(*, G)}. Más simple,
más seguro (no hay fuentes no autorizadas), y usa el rango 232.0.0.0/8.
Requiere IGMPv3 en los hosts.
\end{concepto}

% ─────────────────────────────────────────────
\section{\eh{television} Aplicaciones Reales}

\begin{ejemplo}
\textbf{IPTV (Televisión sobre IP):} operadores como Izzi, Totalplay o Claro TV
distribuyen canales de TV usando multicast. Un canal de HD requiere $\approx$4--8 Mbps.
Con unicast, 10,000 suscriptores viendo el mismo canal = 80 Gbps desde el servidor.
Con multicast = 8 Mbps independientemente del número de suscriptores. El router del
cliente (CPE) hace IGMP Join cuando el usuario cambia de canal.
\end{ejemplo}

\begin{ejemplo}
\textbf{Protocolos de routing:} OSPF usa \texttt{224.0.0.5} (todos los routers OSPF)
y \texttt{224.0.0.6} (DR/BDR). Esto es multicast link-local: el paquete nunca pasa
de ese segmento L2. Sin multicast, OSPF tendría que usar broadcast o unicast punto a punto.
\end{ejemplo}

\begin{moderno}
\emoji{rocket} \textbf{Multicast en datacenters modernos:} VxLAN con BUM traffic
(Broadcast, Unknown-unicast, Multicast) puede usar multicast IP underlay en lugar
de ingress replication. El switch Leaf hace IGMP Join para el grupo multicast del VNI.
Sin embargo, EVPN con ingress replication es hoy más común porque no requiere multicast
en el underlay (muchos clouds no lo soportan). En Kubernetes, \textbf{no hay multicast
nativo en las VPCs de nube pública}: se simula con servicios como AWS SNS (fan-out de
mensajes) o con aplicación-level multicast en frameworks como Hazelcast.
\end{moderno}

\section{\eh{check-mark-button} Resumen del Ecosistema Parte 2}

Al completar estos apuntes has cubierto el stack completo del networking moderno:

\begin{center}
\begin{tabular}{ll}
\toprule
\textbf{Tema} & \textbf{Dónde aplica hoy} \\
\midrule
Wireless empresarial   & Oficinas, campus, IoT industrial \\
Switching avanzado     & Cualquier red LAN empresarial \\
SDN / OpenFlow         & Datacenter, WAN empresarial \\
VxLAN + EVPN           & AWS, Azure, GCP internamente; datacenter privado \\
MPLS + QoS             & ISPs, WAN empresarial, voz sobre IP \\
Kubernetes networking  & Todo cluster en producción \\
Cloud networking       & Toda aplicación en AWS/GCP/Azure \\
Multicast              & IPTV, protocolos de routing, streaming masivo \\
\bottomrule
\end{tabular}
\end{center}


\end{document}
